% ============================================================
%  SECCIÓN: RESULTADOS Y ANÁLISIS — ENSAYO DE CONSUMO DE POTENCIA
%  Tesis: Equipo de Ensayos de Seguridad Eléctrica Portátil
%  Autor: Valentina Ferrini
% ============================================================
%
% Paquetes requeridos en el preámbulo principal:
%   \usepackage{booktabs}
%   \usepackage{multirow}
%   \usepackage{array}
%   \usepackage{float}
%   \usepackage{graphicx}
%   \usepackage{amsmath}
%   \usepackage{amssymb}
%   \usepackage{siunitx}
%   \usepackage[table]{xcolor}
%   \usepackage{pgfplots}        % para gráficos (opcional, ver nota)
%   \pgfplotsset{compat=1.18}
%   \usepackage[spanish]{babel}
%
% INSTRUCCIONES PARA COMPLETAR:
%   Buscar todas las ocurrencias de \PLACEHOLDER{...} y reemplazar
%   con los valores reales obtenidos durante el ensayo.
%   Los comentarios junto a cada placeholder indican el tipo de dato
%   esperado y su unidad.
%
% Uso: % ============================================================
%  SECCIÓN: RESULTADOS Y ANÁLISIS — ENSAYO DE CONSUMO DE POTENCIA
%  Tesis: Equipo de Ensayos de Seguridad Eléctrica Portátil
%  Autor: Valentina Ferrini
% ============================================================
%
% Paquetes requeridos en el preámbulo principal:
%   \usepackage{booktabs}
%   \usepackage{multirow}
%   \usepackage{array}
%   \usepackage{float}
%   \usepackage{graphicx}
%   \usepackage{amsmath}
%   \usepackage{amssymb}
%   \usepackage{siunitx}
%   \usepackage[table]{xcolor}
%   \usepackage{pgfplots}        % para gráficos (opcional, ver nota)
%   \pgfplotsset{compat=1.18}
%   \usepackage[spanish]{babel}
%
% INSTRUCCIONES PARA COMPLETAR:
%   Buscar todas las ocurrencias de \PLACEHOLDER{...} y reemplazar
%   con los valores reales obtenidos durante el ensayo.
%   Los comentarios junto a cada placeholder indican el tipo de dato
%   esperado y su unidad.
%
% Uso: % ============================================================
%  SECCIÓN: RESULTADOS Y ANÁLISIS — ENSAYO DE CONSUMO DE POTENCIA
%  Tesis: Equipo de Ensayos de Seguridad Eléctrica Portátil
%  Autor: Valentina Ferrini
% ============================================================
%
% Paquetes requeridos en el preámbulo principal:
%   \usepackage{booktabs}
%   \usepackage{multirow}
%   \usepackage{array}
%   \usepackage{float}
%   \usepackage{graphicx}
%   \usepackage{amsmath}
%   \usepackage{amssymb}
%   \usepackage{siunitx}
%   \usepackage[table]{xcolor}
%   \usepackage{pgfplots}        % para gráficos (opcional, ver nota)
%   \pgfplotsset{compat=1.18}
%   \usepackage[spanish]{babel}
%
% INSTRUCCIONES PARA COMPLETAR:
%   Buscar todas las ocurrencias de \PLACEHOLDER{...} y reemplazar
%   con los valores reales obtenidos durante el ensayo.
%   Los comentarios junto a cada placeholder indican el tipo de dato
%   esperado y su unidad.
%
% Uso: % ============================================================
%  SECCIÓN: RESULTADOS Y ANÁLISIS — ENSAYO DE CONSUMO DE POTENCIA
%  Tesis: Equipo de Ensayos de Seguridad Eléctrica Portátil
%  Autor: Valentina Ferrini
% ============================================================
%
% Paquetes requeridos en el preámbulo principal:
%   \usepackage{booktabs}
%   \usepackage{multirow}
%   \usepackage{array}
%   \usepackage{float}
%   \usepackage{graphicx}
%   \usepackage{amsmath}
%   \usepackage{amssymb}
%   \usepackage{siunitx}
%   \usepackage[table]{xcolor}
%   \usepackage{pgfplots}        % para gráficos (opcional, ver nota)
%   \pgfplotsset{compat=1.18}
%   \usepackage[spanish]{babel}
%
% INSTRUCCIONES PARA COMPLETAR:
%   Buscar todas las ocurrencias de \PLACEHOLDER{...} y reemplazar
%   con los valores reales obtenidos durante el ensayo.
%   Los comentarios junto a cada placeholder indican el tipo de dato
%   esperado y su unidad.
%
% Uso: \input{resultados_analisis}
% ============================================================

% Comando auxiliar para resaltar placeholders visualmente en Overleaf
% (comentar o eliminar antes de la entrega final)
\newcommand{\PLACEHOLDER}[1]{\textcolor{red}{\textbf{[#1]}}}

% ============================================================
\section{Resultados y Análisis}
\label{sec:resultados}
% ============================================================

Esta sección presenta los resultados obtenidos durante la campaña
de ensayos de validación del prototipo desarrollado.
Se exponen las mediciones realizadas sobre el dispositivo bajo
prueba (DUT), se comparan con los valores de referencia
establecidos por la norma IEC~60335-1 y con un instrumento patrón,
y se analiza el desempeño del sistema en términos de exactitud,
repetibilidad y cumplimiento normativo.

% ------------------------------------------------------------
\subsection{Configuración experimental}
\label{subsec:config_experimental}
% ------------------------------------------------------------

\subsubsection{Dispositivo bajo prueba (DUT)}

Las mediciones se realizaron sobre un electrodoméstico de uso
doméstico de la categoría \textbf{\PLACEHOLDER{nombre del aparato,
p.ej.: plancha a vapor}}, cuyas características declaradas por
el fabricante se resumen en la Tabla~\ref{tab:dut_datos}.

\begin{table}[H]
\centering
\caption{Datos de placa del dispositivo bajo prueba}
\label{tab:dut_datos}
\renewcommand{\arraystretch}{1.35}
\begin{tabular}{@{} l l @{}}
\toprule
\textbf{Parámetro} & \textbf{Valor declarado} \\
\midrule
Marca / modelo          & \PLACEHOLDER{Marca — Modelo} \\
Clase de aislamiento    & \PLACEHOLDER{Clase I / Clase II} \\
Tensión nominal         & \PLACEHOLDER{220 / 230} V AC, 50 Hz \\
Potencia nominal        & \PLACEHOLDER{XXX} W \\
Corriente nominal       & \PLACEHOLDER{X{,}X} A \\
Factor de potencia      & \PLACEHOLDER{X{,}XX} (si figura en placa) \\
Norma de producto       & \PLACEHOLDER{IEC 60335-2-3 (planchas) u otra} \\
\bottomrule
\end{tabular}
\end{table}

\subsubsection{Instrumentación utilizada}

El ensayo se ejecutó con el prototipo desarrollado en este proyecto.
Adicionalmente, como instrumento de referencia para la validación
de las mediciones, se empleó \PLACEHOLDER{nombre y modelo del
instrumento patrón, p.ej.: analizador de potencia Fluke 435-II /
multímetro Fluke 87V / vatímetro patrón del laboratorio}.

Las especificaciones de exactitud del instrumento patrón se
detallan en la Tabla~\ref{tab:instrumento_patron}.

\begin{table}[H]
\centering
\caption{Especificaciones del instrumento de referencia utilizado}
\label{tab:instrumento_patron}
\renewcommand{\arraystretch}{1.35}
\begin{tabular}{@{} l l @{}}
\toprule
\textbf{Parámetro} & \textbf{Especificación} \\
\midrule
Marca / modelo              & \PLACEHOLDER{Marca — Modelo} \\
Exactitud en tensión        & \PLACEHOLDER{$\pm$X{,}X\% + X dígitos} \\
Exactitud en corriente      & \PLACEHOLDER{$\pm$X{,}X\% + X dígitos} \\
Exactitud en potencia       & \PLACEHOLDER{$\pm$X{,}X\%} \\
Trazabilidad de calibración & \PLACEHOLDER{Calibrado por laboratorio
                              acreditado OAA N\textsuperscript{o}~XXX,
                              fecha: MM/AAAA} \\
\bottomrule
\end{tabular}
\end{table}

\subsubsection{Condiciones ambientales}

Las condiciones del laboratorio durante la realización de los
ensayos se registraron al inicio y al final de la sesión de
mediciones y se presentan en la Tabla~\ref{tab:condiciones_amb}.

\begin{table}[H]
\centering
\caption{Condiciones ambientales durante el ensayo}
\label{tab:condiciones_amb}
\renewcommand{\arraystretch}{1.35}
\begin{tabular}{@{} l r r r @{}}
\toprule
\textbf{Parámetro} & \textbf{Inicio} & \textbf{Final} & \textbf{Límite IEC} \\
\midrule
Temperatura ambiente [°C]   & \PLACEHOLDER{XX,X} & \PLACEHOLDER{XX,X} & $(20 \pm 5)$~°C \\
Humedad relativa [\%]        & \PLACEHOLDER{XX}   & \PLACEHOLDER{XX}   & $< 85$\% \\
Tensión de red [V$_{\text{rms}}$] & \PLACEHOLDER{XXX,X} & \PLACEHOLDER{XXX,X} & $V_n \pm 1\%$ \\
Frecuencia de red [Hz]       & \PLACEHOLDER{XX,X} & \PLACEHOLDER{XX,X} & $50 \pm 0{,}5$~Hz \\
\bottomrule
\end{tabular}
\end{table}

% ------------------------------------------------------------
\subsection{Resultados del ensayo de consumo de potencia}
\label{subsec:resultados_potencia}
% ------------------------------------------------------------

\subsubsection{Mediciones individuales}

Se realizaron \PLACEHOLDER{N} mediciones consecutivas
con el DUT en condición térmica estable
(temperatura de operación normal alcanzada tras
\PLACEHOLDER{X}~minutos de precalentamiento).
Cada medición corresponde a un ciclo completo de adquisición
del módulo PZEM-004T~v3.0, con un intervalo de muestreo de
\PLACEHOLDER{X}~segundos.
Los valores registrados se presentan en la
Tabla~\ref{tab:mediciones_pzem}.

\begin{table}[H]
\centering
\caption{Mediciones de parámetros eléctricos — DUT: \PLACEHOLDER{nombre del aparato}}
\label{tab:mediciones_pzem}
\renewcommand{\arraystretch}{1.3}
\rowcolors{2}{gray!8}{white}
\begin{tabular}{@{} c r r r r r @{}}
\toprule
\textbf{Muestra} &
\textbf{$V_{\text{rms}}$ [V]} &
\textbf{$I_{\text{rms}}$ [A]} &
\textbf{$P$ [W]} &
\textbf{$\cos\varphi$} &
\textbf{$f$ [Hz]} \\
\midrule
1  & \PLACEHOLDER{XXX,X} & \PLACEHOLDER{X,XXX} & \PLACEHOLDER{XXXX,X} & \PLACEHOLDER{X,XX} & \PLACEHOLDER{XX,X} \\
2  & \PLACEHOLDER{XXX,X} & \PLACEHOLDER{X,XXX} & \PLACEHOLDER{XXXX,X} & \PLACEHOLDER{X,XX} & \PLACEHOLDER{XX,X} \\
3  & \PLACEHOLDER{XXX,X} & \PLACEHOLDER{X,XXX} & \PLACEHOLDER{XXXX,X} & \PLACEHOLDER{X,XX} & \PLACEHOLDER{XX,X} \\
4  & \PLACEHOLDER{XXX,X} & \PLACEHOLDER{X,XXX} & \PLACEHOLDER{XXXX,X} & \PLACEHOLDER{X,XX} & \PLACEHOLDER{XX,X} \\
5  & \PLACEHOLDER{XXX,X} & \PLACEHOLDER{X,XXX} & \PLACEHOLDER{XXXX,X} & \PLACEHOLDER{X,XX} & \PLACEHOLDER{XX,X} \\
6  & \PLACEHOLDER{XXX,X} & \PLACEHOLDER{X,XXX} & \PLACEHOLDER{XXXX,X} & \PLACEHOLDER{X,XX} & \PLACEHOLDER{XX,X} \\
7  & \PLACEHOLDER{XXX,X} & \PLACEHOLDER{X,XXX} & \PLACEHOLDER{XXXX,X} & \PLACEHOLDER{X,XX} & \PLACEHOLDER{XX,X} \\
8  & \PLACEHOLDER{XXX,X} & \PLACEHOLDER{X,XXX} & \PLACEHOLDER{XXXX,X} & \PLACEHOLDER{X,XX} & \PLACEHOLDER{XX,X} \\
9  & \PLACEHOLDER{XXX,X} & \PLACEHOLDER{X,XXX} & \PLACEHOLDER{XXXX,X} & \PLACEHOLDER{X,XX} & \PLACEHOLDER{XX,X} \\
10 & \PLACEHOLDER{XXX,X} & \PLACEHOLDER{X,XXX} & \PLACEHOLDER{XXXX,X} & \PLACEHOLDER{X,XX} & \PLACEHOLDER{XX,X} \\
\midrule
\textbf{Media} $\bar{x}$
   & \PLACEHOLDER{XXX,X} & \PLACEHOLDER{X,XXX} & \PLACEHOLDER{XXXX,X} & \PLACEHOLDER{X,XX} & \PLACEHOLDER{XX,X} \\
\textbf{Desv. est.} $s$
   & \PLACEHOLDER{X,X}   & \PLACEHOLDER{X,XXX} & \PLACEHOLDER{X,X}    & \PLACEHOLDER{X,XX} & \PLACEHOLDER{X,X}  \\
\textbf{Mín.}
   & \PLACEHOLDER{XXX,X} & \PLACEHOLDER{X,XXX} & \PLACEHOLDER{XXXX,X} & \PLACEHOLDER{X,XX} & \PLACEHOLDER{XX,X} \\
\textbf{Máx.}
   & \PLACEHOLDER{XXX,X} & \PLACEHOLDER{X,XXX} & \PLACEHOLDER{XXXX,X} & \PLACEHOLDER{X,XX} & \PLACEHOLDER{XX,X} \\
\bottomrule
\end{tabular}
\end{table}

\subsubsection{Análisis estadístico}

A partir de las \PLACEHOLDER{N} mediciones registradas,
se calcularon los estadísticos descriptivos de la potencia activa $P$:

\begin{align}
    \bar{P} &= \frac{1}{N}\sum_{i=1}^{N} P_i
             = \PLACEHOLDER{XXXX{,}X}~\text{W}
    \label{eq:media_potencia} \\[6pt]
    s_P &= \sqrt{\frac{1}{N-1}\sum_{i=1}^{N}(P_i - \bar{P})^2}
          = \PLACEHOLDER{X{,}X}~\text{W}
    \label{eq:std_potencia} \\[6pt]
    u_{\bar{P}} &= \frac{s_P}{\sqrt{N}}
                 = \PLACEHOLDER{X{,}X}~\text{W}
    \label{eq:incertidumbre_estadistica}
\end{align}

\noindent donde $u_{\bar{P}}$ es la incertidumbre estándar tipo~A
asociada a la dispersión de las mediciones.

El \textbf{coeficiente de variación} (CV) de la potencia medida es:

\begin{equation}
    \text{CV} = \frac{s_P}{\bar{P}} \times 100\%
              = \PLACEHOLDER{X{,}X}\%
    \label{eq:cv_potencia}
\end{equation}

Un CV menor al 2\% indica una buena repetibilidad del sistema
de medición bajo las condiciones de ensayo.

% ------------------------------------------------------------
\subsection{Validación frente al instrumento de referencia}
\label{subsec:validacion_patron}
% ------------------------------------------------------------

\subsubsection{Comparación de mediciones}

Con el objetivo de verificar la exactitud del prototipo,
se realizaron mediciones simultáneas con el equipo desarrollado
y con el instrumento de referencia sobre el mismo DUT y en las
mismas condiciones de carga.
Los resultados comparativos se presentan en la
Tabla~\ref{tab:comparacion_patron}.

\begin{table}[H]
\centering
\caption{Comparación entre el prototipo y el instrumento de referencia}
\label{tab:comparacion_patron}
\renewcommand{\arraystretch}{1.4}
\begin{tabular}{@{} l r r r r @{}}
\toprule
\textbf{Parámetro}
    & \textbf{Prototipo}
    & \textbf{Referencia}
    & \textbf{Error abs.}
    & \textbf{Error rel. [\%]} \\
\midrule
$V_{\text{rms}}$ [V]
    & \PLACEHOLDER{XXX,X}
    & \PLACEHOLDER{XXX,X}
    & \PLACEHOLDER{$\pm$X,X}
    & \PLACEHOLDER{X,XX} \\
$I_{\text{rms}}$ [A]
    & \PLACEHOLDER{X,XXX}
    & \PLACEHOLDER{X,XXX}
    & \PLACEHOLDER{$\pm$X,XXX}
    & \PLACEHOLDER{X,XX} \\
$P$ [W]
    & \PLACEHOLDER{XXXX,X}
    & \PLACEHOLDER{XXXX,X}
    & \PLACEHOLDER{$\pm$X,X}
    & \PLACEHOLDER{X,XX} \\
$\cos\varphi$
    & \PLACEHOLDER{X,XX}
    & \PLACEHOLDER{X,XX}
    & \PLACEHOLDER{$\pm$X,XX}
    & \PLACEHOLDER{X,XX} \\
\bottomrule
\end{tabular}
\end{table}

\noindent El error relativo en potencia activa $P$ se calcula como:

\begin{equation}
    e_r = \frac{P_{\text{prototipo}} - P_{\text{referencia}}}
               {P_{\text{referencia}}} \times 100\%
        = \PLACEHOLDER{X{,}XX}\%
    \label{eq:error_relativo}
\end{equation}

La especificación de exactitud del módulo PZEM-004T~v3.0 en
potencia activa es de $\pm$0{,}5\%.
El error relativo obtenido de \PLACEHOLDER{X{,}XX}\%
\PLACEHOLDER{es / no es} consistente con dicha especificación,
\PLACEHOLDER{lo que confirma / lo que sugiere una posible causa
de error en: descripción del problema observado}.

\subsubsection{Cumplimiento del criterio normativo}

La norma IEC~60335-1 (cláusula~10.2) establece que la potencia
consumida por el DUT no debe diferir de la potencia declarada
por el fabricante en más de $\pm$10\%.
La evaluación de este criterio se resume en la
Tabla~\ref{tab:criterio_normativo}.

\begin{table}[H]
\centering
\caption{Evaluación del criterio normativo de consumo — IEC~60335-1, cláusula~10.2}
\label{tab:criterio_normativo}
\renewcommand{\arraystretch}{1.4}
\begin{tabular}{@{} l r @{}}
\toprule
\textbf{Parámetro} & \textbf{Valor} \\
\midrule
Potencia declarada por el fabricante ($P_{\text{dec}}$)
    & \PLACEHOLDER{XXXX} W \\
Límite inferior admisible ($P_{\text{dec}} \times 0{,}90$)
    & \PLACEHOLDER{XXXX} W \\
Límite superior admisible ($P_{\text{dec}} \times 1{,}10$)
    & \PLACEHOLDER{XXXX} W \\
Potencia media medida por el prototipo ($\bar{P}$)
    & \PLACEHOLDER{XXXX,X} W \\
Desviación respecto al valor declarado
    & \PLACEHOLDER{$+$X{,}X / $-$X{,}X}\% \\
\midrule
\textbf{Resultado del ensayo}
    & \PLACEHOLDER{\textcolor{green!60!black}{\textbf{APROBADO}} /
                   \textcolor{red}{\textbf{REPROBADO}}} \\
\bottomrule
\end{tabular}
\end{table}

% ------------------------------------------------------------
\subsection{Análisis de incertidumbre de medición}
\label{subsec:incertidumbre}
% ------------------------------------------------------------

El análisis de incertidumbre se realiza siguiendo el marco de la
Guía para la Expresión de la Incertidumbre de Medición
(GUM, JCGM~100:2008), considerando las fuentes de incertidumbre
relevantes para la medición de potencia activa.

\subsubsection{Fuentes de incertidumbre identificadas}

\begin{itemize}
    \item \textbf{Tipo~A — dispersión estadística de las mediciones:}
          incertidumbre debida a la variabilidad aleatoria del proceso
          de medición, cuantificada por la desviación estándar de la
          media $u_{\bar{P}}$ (ecuación~\ref{eq:incertidumbre_estadistica}).

    \item \textbf{Tipo~B — exactitud del sensor PZEM-004T:}
          la hoja de datos del módulo especifica una exactitud de
          $\pm$0{,}5\% en potencia activa.
          Asumiendo distribución rectangular, la incertidumbre
          estándar asociada es:
          \begin{equation}
              u_B = \frac{0{,}005 \times \bar{P}}{\sqrt{3}}
                  = \PLACEHOLDER{X{,}X}~\text{W}
              \label{eq:incertidumbre_B}
          \end{equation}

    \item \textbf{Tipo~B — variación de la tensión de red:}
          fluctuaciones de la tensión de alimentación durante el
          ensayo se propagan a la medición de potencia.
          Con una variación observada de $\Delta V = \PLACEHOLDER{X{,}X}$~V
          y el coeficiente de sensibilidad
          $c_V = \partial P / \partial V \approx I_{\text{rms}}$:
          \begin{equation}
              u_{V} = c_V \cdot \frac{\Delta V}{\sqrt{3}}
                    = \PLACEHOLDER{X{,}X}~\text{W}
              \label{eq:incertidumbre_V}
          \end{equation}

    \item \textbf{Tipo~B — resolución del sensor:}
          la resolución en potencia del PZEM-004T es 0{,}1~W.
          Asumiendo distribución rectangular:
          \begin{equation}
              u_{\text{res}} = \frac{0{,}1}{2\sqrt{3}}
                             \approx 0{,}029~\text{W}
              \label{eq:incertidumbre_res}
          \end{equation}
\end{itemize}

\subsubsection{Incertidumbre combinada y expandida}

Considerando las fuentes anteriores como independientes entre sí,
la incertidumbre estándar combinada es:

\begin{equation}
    u_c = \sqrt{u_{\bar{P}}^2 + u_B^2 + u_V^2 + u_{\text{res}}^2}
        = \PLACEHOLDER{X{,}X}~\text{W}
    \label{eq:incertidumbre_combinada}
\end{equation}

La incertidumbre expandida, con factor de cobertura $k=2$
(nivel de confianza $\approx$95\%, distribución normal), es:

\begin{equation}
    U = k \cdot u_c = 2 \times \PLACEHOLDER{X{,}X}
      = \PLACEHOLDER{X{,}X}~\text{W}
    \label{eq:incertidumbre_expandida}
\end{equation}

El resultado final de la medición de potencia se expresa como:

\begin{equation}
    P = (\PLACEHOLDER{XXXX{,}X} \pm \PLACEHOLDER{X{,}X})~\text{W}
        \quad (k=2,\; p \approx 95\%)
    \label{eq:resultado_final}
\end{equation}

\subsubsection{Relación entre incertidumbre y límite normativo}

La incertidumbre expandida del sistema ($U = \PLACEHOLDER{X{,}X}$~W)
representa el \PLACEHOLDER{X{,}X}\% del límite normativo
($\pm$10\% de $P_{\text{dec}} = \PLACEHOLDER{XXXX}$~W).
Este valor es \PLACEHOLDER{aceptable para los propósitos del
prototipo / recomendable reducir mediante calibración del sensor},
dado que el criterio general de metrología establece que la
incertidumbre del instrumento de medición debe ser al menos
cuatro veces menor que la tolerancia que se desea verificar.

% ------------------------------------------------------------
\subsection{Evaluación de la repetibilidad del sistema}
\label{subsec:repetibilidad}
% ------------------------------------------------------------

La repetibilidad del sistema se evalúa como la capacidad de
obtener resultados consistentes bajo las mismas condiciones
experimentales en mediciones sucesivas.

\begin{table}[H]
\centering
\caption{Indicadores de repetibilidad del ensayo de potencia}
\label{tab:repetibilidad}
\renewcommand{\arraystretch}{1.35}
\begin{tabular}{@{} l r l @{}}
\toprule
\textbf{Indicador} & \textbf{Valor} & \textbf{Interpretación} \\
\midrule
Media $\bar{P}$                 & \PLACEHOLDER{XXXX,X} W  & Valor central de la distribución \\
Desviación estándar $s_P$       & \PLACEHOLDER{X,X} W     & Dispersión de las mediciones \\
Coeficiente de variación CV     & \PLACEHOLDER{X,X}\%     & Repetibilidad relativa \\
Rango ($P_{\max} - P_{\min}$)   & \PLACEHOLDER{X,X} W     & Amplitud total de dispersión \\
Incertidumbre estándar $u_{\bar{P}}$ & \PLACEHOLDER{X,X} W & Incertidumbre tipo~A \\
\midrule
\textbf{Criterio de repetibilidad} & CV $<$ 2\%
    & \PLACEHOLDER{\textcolor{green!60!black}{\textbf{CUMPLE}} /
                   \textcolor{red}{\textbf{NO CUMPLE}}} \\
\bottomrule
\end{tabular}
\end{table}

% ------------------------------------------------------------
\subsection{Discusión de resultados}
\label{subsec:discusion}
% ------------------------------------------------------------

Los resultados obtenidos permiten extraer las siguientes
conclusiones sobre el desempeño del prototipo:

\paragraph{Exactitud en la medición de potencia.}
El error relativo promedio frente al instrumento de referencia
fue de \PLACEHOLDER{X{,}XX}\%, valor que
\PLACEHOLDER{se encuentra dentro / supera levemente}
la especificación del fabricante del módulo PZEM-004T (±0{,}5\%).
\PLACEHOLDER{Describir aquí la causa probable del error observado
si supera la especificación, p.ej.: caída de tensión en el cableado,
temperatura de operación del módulo, saturación del TC, etc.}

\paragraph{Repetibilidad.}
El coeficiente de variación obtenido de \PLACEHOLDER{X{,}X}\%
\PLACEHOLDER{confirma / no confirma} una repetibilidad adecuada
del sistema para el propósito del ensayo.
\PLACEHOLDER{Si CV > 2\%: identificar causa de la variabilidad
(fluctuaciones de red, termostato del DUT activo durante la medición,
interferencias en el bus UART, etc.) y proponer acción correctiva.}

\paragraph{Cumplimiento normativo.}
El DUT ensayado \PLACEHOLDER{aprobó / reprobó} el criterio de
consumo de la norma IEC~60335-1 (cláusula~10.2), con una desviación
de \PLACEHOLDER{X{,}X}\% respecto al valor declarado,
dentro del límite admisible de $\pm$10\%.
Este resultado \PLACEHOLDER{es / no es} consistente con lo esperado
para un electrodoméstico \PLACEHOLDER{en buenas condiciones / con
el defecto identificado: descripción}.

\paragraph{Limitaciones identificadas.}
Durante la campaña de ensayos se identificaron las siguientes
limitaciones del prototipo en su estado actual:

\begin{itemize}
    \item \PLACEHOLDER{Limitación 1: p.ej., el módulo PZEM-004T
          presenta lecturas inestables durante los primeros 30~s
          de conexión, por lo que se implementó un retardo de
          estabilización en el software.}
    \item \PLACEHOLDER{Limitación 2: p.ej., la precisión de la
          medición de factor de potencia se degrada para cargas
          muy resistivas ($\cos\varphi > 0{,}98$) debido a la
          resolución de 0{,}01 del módulo.}
    \item \PLACEHOLDER{Limitación 3: agregar según observaciones
          reales del ensayo.}
\end{itemize}

Estas limitaciones no comprometen la validez del ensayo de consumo
como función principal del prototipo, pero deben ser consideradas
al momento de ampliar el sistema hacia los ensayos de fuga y
rigidez dieléctrica contemplados en la arquitectura del equipo.

% ------------------------------------------------------------
\subsection{Comparación con el criterio de éxito del proyecto}
\label{subsec:criterio_exito}
% ------------------------------------------------------------

La Tabla~\ref{tab:criterios_exito} evalúa los resultados obtenidos
frente a los criterios de éxito definidos en la
Sección~\ref{sec:criterios_evaluacion} del presente informe.

\begin{table}[H]
\centering
\caption{Evaluación de criterios de éxito del prototipo}
\label{tab:criterios_exito}
\renewcommand{\arraystretch}{1.4}
\begin{tabular}{@{} p{5cm} p{3cm} p{3cm} p{1.8cm} @{}}
\toprule
\textbf{Criterio} & \textbf{Meta definida} & \textbf{Resultado obtenido} & \textbf{Estado} \\
\midrule
Reducción del tiempo de prueba vs.\ método manual
    & $\geq 30\%$
    & \PLACEHOLDER{XX\%}
    & \PLACEHOLDER{\textcolor{green!60!black}{\checkmark} /
                   \textcolor{red}{$\times$}} \\
\addlinespace
Reducción de intervención del operario
    & $\geq 80\%$
    & \PLACEHOLDER{XX\%}
    & \PLACEHOLDER{\textcolor{green!60!black}{\checkmark} /
                   \textcolor{red}{$\times$}} \\
\addlinespace
Error relativo en potencia frente a patrón
    & Según especif.\ PZEM ($\pm$0{,}5\%)
    & \PLACEHOLDER{$\pm$X{,}XX\%}
    & \PLACEHOLDER{\textcolor{green!60!black}{\checkmark} /
                   \textcolor{red}{$\times$}} \\
\addlinespace
Cumplimiento normativo (criterio IEC~10.2)
    & Detección correcta APROBADO / REPROBADO
    & \PLACEHOLDER{Correcto / Incorrecto}
    & \PLACEHOLDER{\textcolor{green!60!black}{\checkmark} /
                   \textcolor{red}{$\times$}} \\
\addlinespace
Repetibilidad (CV)
    & CV $< 2\%$
    & \PLACEHOLDER{X{,}X\%}
    & \PLACEHOLDER{\textcolor{green!60!black}{\checkmark} /
                   \textcolor{red}{$\times$}} \\
\addlinespace
Ajuste al presupuesto
    & $\leq$ USD~1030
    & USD~\PLACEHOLDER{XXX}
    & \PLACEHOLDER{\textcolor{green!60!black}{\checkmark} /
                   \textcolor{red}{$\times$}} \\
\bottomrule
\end{tabular}
\end{table}

% ============================================================

% Comando auxiliar para resaltar placeholders visualmente en Overleaf
% (comentar o eliminar antes de la entrega final)
\newcommand{\PLACEHOLDER}[1]{\textcolor{red}{\textbf{[#1]}}}

% ============================================================
\section{Resultados y Análisis}
\label{sec:resultados}
% ============================================================

Esta sección presenta los resultados obtenidos durante la campaña
de ensayos de validación del prototipo desarrollado.
Se exponen las mediciones realizadas sobre el dispositivo bajo
prueba (DUT), se comparan con los valores de referencia
establecidos por la norma IEC~60335-1 y con un instrumento patrón,
y se analiza el desempeño del sistema en términos de exactitud,
repetibilidad y cumplimiento normativo.

% ------------------------------------------------------------
\subsection{Configuración experimental}
\label{subsec:config_experimental}
% ------------------------------------------------------------

\subsubsection{Dispositivo bajo prueba (DUT)}

Las mediciones se realizaron sobre un electrodoméstico de uso
doméstico de la categoría \textbf{\PLACEHOLDER{nombre del aparato,
p.ej.: plancha a vapor}}, cuyas características declaradas por
el fabricante se resumen en la Tabla~\ref{tab:dut_datos}.

\begin{table}[H]
\centering
\caption{Datos de placa del dispositivo bajo prueba}
\label{tab:dut_datos}
\renewcommand{\arraystretch}{1.35}
\begin{tabular}{@{} l l @{}}
\toprule
\textbf{Parámetro} & \textbf{Valor declarado} \\
\midrule
Marca / modelo          & \PLACEHOLDER{Marca — Modelo} \\
Clase de aislamiento    & \PLACEHOLDER{Clase I / Clase II} \\
Tensión nominal         & \PLACEHOLDER{220 / 230} V AC, 50 Hz \\
Potencia nominal        & \PLACEHOLDER{XXX} W \\
Corriente nominal       & \PLACEHOLDER{X{,}X} A \\
Factor de potencia      & \PLACEHOLDER{X{,}XX} (si figura en placa) \\
Norma de producto       & \PLACEHOLDER{IEC 60335-2-3 (planchas) u otra} \\
\bottomrule
\end{tabular}
\end{table}

\subsubsection{Instrumentación utilizada}

El ensayo se ejecutó con el prototipo desarrollado en este proyecto.
Adicionalmente, como instrumento de referencia para la validación
de las mediciones, se empleó \PLACEHOLDER{nombre y modelo del
instrumento patrón, p.ej.: analizador de potencia Fluke 435-II /
multímetro Fluke 87V / vatímetro patrón del laboratorio}.

Las especificaciones de exactitud del instrumento patrón se
detallan en la Tabla~\ref{tab:instrumento_patron}.

\begin{table}[H]
\centering
\caption{Especificaciones del instrumento de referencia utilizado}
\label{tab:instrumento_patron}
\renewcommand{\arraystretch}{1.35}
\begin{tabular}{@{} l l @{}}
\toprule
\textbf{Parámetro} & \textbf{Especificación} \\
\midrule
Marca / modelo              & \PLACEHOLDER{Marca — Modelo} \\
Exactitud en tensión        & \PLACEHOLDER{$\pm$X{,}X\% + X dígitos} \\
Exactitud en corriente      & \PLACEHOLDER{$\pm$X{,}X\% + X dígitos} \\
Exactitud en potencia       & \PLACEHOLDER{$\pm$X{,}X\%} \\
Trazabilidad de calibración & \PLACEHOLDER{Calibrado por laboratorio
                              acreditado OAA N\textsuperscript{o}~XXX,
                              fecha: MM/AAAA} \\
\bottomrule
\end{tabular}
\end{table}

\subsubsection{Condiciones ambientales}

Las condiciones del laboratorio durante la realización de los
ensayos se registraron al inicio y al final de la sesión de
mediciones y se presentan en la Tabla~\ref{tab:condiciones_amb}.

\begin{table}[H]
\centering
\caption{Condiciones ambientales durante el ensayo}
\label{tab:condiciones_amb}
\renewcommand{\arraystretch}{1.35}
\begin{tabular}{@{} l r r r @{}}
\toprule
\textbf{Parámetro} & \textbf{Inicio} & \textbf{Final} & \textbf{Límite IEC} \\
\midrule
Temperatura ambiente [°C]   & \PLACEHOLDER{XX,X} & \PLACEHOLDER{XX,X} & $(20 \pm 5)$~°C \\
Humedad relativa [\%]        & \PLACEHOLDER{XX}   & \PLACEHOLDER{XX}   & $< 85$\% \\
Tensión de red [V$_{\text{rms}}$] & \PLACEHOLDER{XXX,X} & \PLACEHOLDER{XXX,X} & $V_n \pm 1\%$ \\
Frecuencia de red [Hz]       & \PLACEHOLDER{XX,X} & \PLACEHOLDER{XX,X} & $50 \pm 0{,}5$~Hz \\
\bottomrule
\end{tabular}
\end{table}

% ------------------------------------------------------------
\subsection{Resultados del ensayo de consumo de potencia}
\label{subsec:resultados_potencia}
% ------------------------------------------------------------

\subsubsection{Mediciones individuales}

Se realizaron \PLACEHOLDER{N} mediciones consecutivas
con el DUT en condición térmica estable
(temperatura de operación normal alcanzada tras
\PLACEHOLDER{X}~minutos de precalentamiento).
Cada medición corresponde a un ciclo completo de adquisición
del módulo PZEM-004T~v3.0, con un intervalo de muestreo de
\PLACEHOLDER{X}~segundos.
Los valores registrados se presentan en la
Tabla~\ref{tab:mediciones_pzem}.

\begin{table}[H]
\centering
\caption{Mediciones de parámetros eléctricos — DUT: \PLACEHOLDER{nombre del aparato}}
\label{tab:mediciones_pzem}
\renewcommand{\arraystretch}{1.3}
\rowcolors{2}{gray!8}{white}
\begin{tabular}{@{} c r r r r r @{}}
\toprule
\textbf{Muestra} &
\textbf{$V_{\text{rms}}$ [V]} &
\textbf{$I_{\text{rms}}$ [A]} &
\textbf{$P$ [W]} &
\textbf{$\cos\varphi$} &
\textbf{$f$ [Hz]} \\
\midrule
1  & \PLACEHOLDER{XXX,X} & \PLACEHOLDER{X,XXX} & \PLACEHOLDER{XXXX,X} & \PLACEHOLDER{X,XX} & \PLACEHOLDER{XX,X} \\
2  & \PLACEHOLDER{XXX,X} & \PLACEHOLDER{X,XXX} & \PLACEHOLDER{XXXX,X} & \PLACEHOLDER{X,XX} & \PLACEHOLDER{XX,X} \\
3  & \PLACEHOLDER{XXX,X} & \PLACEHOLDER{X,XXX} & \PLACEHOLDER{XXXX,X} & \PLACEHOLDER{X,XX} & \PLACEHOLDER{XX,X} \\
4  & \PLACEHOLDER{XXX,X} & \PLACEHOLDER{X,XXX} & \PLACEHOLDER{XXXX,X} & \PLACEHOLDER{X,XX} & \PLACEHOLDER{XX,X} \\
5  & \PLACEHOLDER{XXX,X} & \PLACEHOLDER{X,XXX} & \PLACEHOLDER{XXXX,X} & \PLACEHOLDER{X,XX} & \PLACEHOLDER{XX,X} \\
6  & \PLACEHOLDER{XXX,X} & \PLACEHOLDER{X,XXX} & \PLACEHOLDER{XXXX,X} & \PLACEHOLDER{X,XX} & \PLACEHOLDER{XX,X} \\
7  & \PLACEHOLDER{XXX,X} & \PLACEHOLDER{X,XXX} & \PLACEHOLDER{XXXX,X} & \PLACEHOLDER{X,XX} & \PLACEHOLDER{XX,X} \\
8  & \PLACEHOLDER{XXX,X} & \PLACEHOLDER{X,XXX} & \PLACEHOLDER{XXXX,X} & \PLACEHOLDER{X,XX} & \PLACEHOLDER{XX,X} \\
9  & \PLACEHOLDER{XXX,X} & \PLACEHOLDER{X,XXX} & \PLACEHOLDER{XXXX,X} & \PLACEHOLDER{X,XX} & \PLACEHOLDER{XX,X} \\
10 & \PLACEHOLDER{XXX,X} & \PLACEHOLDER{X,XXX} & \PLACEHOLDER{XXXX,X} & \PLACEHOLDER{X,XX} & \PLACEHOLDER{XX,X} \\
\midrule
\textbf{Media} $\bar{x}$
   & \PLACEHOLDER{XXX,X} & \PLACEHOLDER{X,XXX} & \PLACEHOLDER{XXXX,X} & \PLACEHOLDER{X,XX} & \PLACEHOLDER{XX,X} \\
\textbf{Desv. est.} $s$
   & \PLACEHOLDER{X,X}   & \PLACEHOLDER{X,XXX} & \PLACEHOLDER{X,X}    & \PLACEHOLDER{X,XX} & \PLACEHOLDER{X,X}  \\
\textbf{Mín.}
   & \PLACEHOLDER{XXX,X} & \PLACEHOLDER{X,XXX} & \PLACEHOLDER{XXXX,X} & \PLACEHOLDER{X,XX} & \PLACEHOLDER{XX,X} \\
\textbf{Máx.}
   & \PLACEHOLDER{XXX,X} & \PLACEHOLDER{X,XXX} & \PLACEHOLDER{XXXX,X} & \PLACEHOLDER{X,XX} & \PLACEHOLDER{XX,X} \\
\bottomrule
\end{tabular}
\end{table}

\subsubsection{Análisis estadístico}

A partir de las \PLACEHOLDER{N} mediciones registradas,
se calcularon los estadísticos descriptivos de la potencia activa $P$:

\begin{align}
    \bar{P} &= \frac{1}{N}\sum_{i=1}^{N} P_i
             = \PLACEHOLDER{XXXX{,}X}~\text{W}
    \label{eq:media_potencia} \\[6pt]
    s_P &= \sqrt{\frac{1}{N-1}\sum_{i=1}^{N}(P_i - \bar{P})^2}
          = \PLACEHOLDER{X{,}X}~\text{W}
    \label{eq:std_potencia} \\[6pt]
    u_{\bar{P}} &= \frac{s_P}{\sqrt{N}}
                 = \PLACEHOLDER{X{,}X}~\text{W}
    \label{eq:incertidumbre_estadistica}
\end{align}

\noindent donde $u_{\bar{P}}$ es la incertidumbre estándar tipo~A
asociada a la dispersión de las mediciones.

El \textbf{coeficiente de variación} (CV) de la potencia medida es:

\begin{equation}
    \text{CV} = \frac{s_P}{\bar{P}} \times 100\%
              = \PLACEHOLDER{X{,}X}\%
    \label{eq:cv_potencia}
\end{equation}

Un CV menor al 2\% indica una buena repetibilidad del sistema
de medición bajo las condiciones de ensayo.

% ------------------------------------------------------------
\subsection{Validación frente al instrumento de referencia}
\label{subsec:validacion_patron}
% ------------------------------------------------------------

\subsubsection{Comparación de mediciones}

Con el objetivo de verificar la exactitud del prototipo,
se realizaron mediciones simultáneas con el equipo desarrollado
y con el instrumento de referencia sobre el mismo DUT y en las
mismas condiciones de carga.
Los resultados comparativos se presentan en la
Tabla~\ref{tab:comparacion_patron}.

\begin{table}[H]
\centering
\caption{Comparación entre el prototipo y el instrumento de referencia}
\label{tab:comparacion_patron}
\renewcommand{\arraystretch}{1.4}
\begin{tabular}{@{} l r r r r @{}}
\toprule
\textbf{Parámetro}
    & \textbf{Prototipo}
    & \textbf{Referencia}
    & \textbf{Error abs.}
    & \textbf{Error rel. [\%]} \\
\midrule
$V_{\text{rms}}$ [V]
    & \PLACEHOLDER{XXX,X}
    & \PLACEHOLDER{XXX,X}
    & \PLACEHOLDER{$\pm$X,X}
    & \PLACEHOLDER{X,XX} \\
$I_{\text{rms}}$ [A]
    & \PLACEHOLDER{X,XXX}
    & \PLACEHOLDER{X,XXX}
    & \PLACEHOLDER{$\pm$X,XXX}
    & \PLACEHOLDER{X,XX} \\
$P$ [W]
    & \PLACEHOLDER{XXXX,X}
    & \PLACEHOLDER{XXXX,X}
    & \PLACEHOLDER{$\pm$X,X}
    & \PLACEHOLDER{X,XX} \\
$\cos\varphi$
    & \PLACEHOLDER{X,XX}
    & \PLACEHOLDER{X,XX}
    & \PLACEHOLDER{$\pm$X,XX}
    & \PLACEHOLDER{X,XX} \\
\bottomrule
\end{tabular}
\end{table}

\noindent El error relativo en potencia activa $P$ se calcula como:

\begin{equation}
    e_r = \frac{P_{\text{prototipo}} - P_{\text{referencia}}}
               {P_{\text{referencia}}} \times 100\%
        = \PLACEHOLDER{X{,}XX}\%
    \label{eq:error_relativo}
\end{equation}

La especificación de exactitud del módulo PZEM-004T~v3.0 en
potencia activa es de $\pm$0{,}5\%.
El error relativo obtenido de \PLACEHOLDER{X{,}XX}\%
\PLACEHOLDER{es / no es} consistente con dicha especificación,
\PLACEHOLDER{lo que confirma / lo que sugiere una posible causa
de error en: descripción del problema observado}.

\subsubsection{Cumplimiento del criterio normativo}

La norma IEC~60335-1 (cláusula~10.2) establece que la potencia
consumida por el DUT no debe diferir de la potencia declarada
por el fabricante en más de $\pm$10\%.
La evaluación de este criterio se resume en la
Tabla~\ref{tab:criterio_normativo}.

\begin{table}[H]
\centering
\caption{Evaluación del criterio normativo de consumo — IEC~60335-1, cláusula~10.2}
\label{tab:criterio_normativo}
\renewcommand{\arraystretch}{1.4}
\begin{tabular}{@{} l r @{}}
\toprule
\textbf{Parámetro} & \textbf{Valor} \\
\midrule
Potencia declarada por el fabricante ($P_{\text{dec}}$)
    & \PLACEHOLDER{XXXX} W \\
Límite inferior admisible ($P_{\text{dec}} \times 0{,}90$)
    & \PLACEHOLDER{XXXX} W \\
Límite superior admisible ($P_{\text{dec}} \times 1{,}10$)
    & \PLACEHOLDER{XXXX} W \\
Potencia media medida por el prototipo ($\bar{P}$)
    & \PLACEHOLDER{XXXX,X} W \\
Desviación respecto al valor declarado
    & \PLACEHOLDER{$+$X{,}X / $-$X{,}X}\% \\
\midrule
\textbf{Resultado del ensayo}
    & \PLACEHOLDER{\textcolor{green!60!black}{\textbf{APROBADO}} /
                   \textcolor{red}{\textbf{REPROBADO}}} \\
\bottomrule
\end{tabular}
\end{table}

% ------------------------------------------------------------
\subsection{Análisis de incertidumbre de medición}
\label{subsec:incertidumbre}
% ------------------------------------------------------------

El análisis de incertidumbre se realiza siguiendo el marco de la
Guía para la Expresión de la Incertidumbre de Medición
(GUM, JCGM~100:2008), considerando las fuentes de incertidumbre
relevantes para la medición de potencia activa.

\subsubsection{Fuentes de incertidumbre identificadas}

\begin{itemize}
    \item \textbf{Tipo~A — dispersión estadística de las mediciones:}
          incertidumbre debida a la variabilidad aleatoria del proceso
          de medición, cuantificada por la desviación estándar de la
          media $u_{\bar{P}}$ (ecuación~\ref{eq:incertidumbre_estadistica}).

    \item \textbf{Tipo~B — exactitud del sensor PZEM-004T:}
          la hoja de datos del módulo especifica una exactitud de
          $\pm$0{,}5\% en potencia activa.
          Asumiendo distribución rectangular, la incertidumbre
          estándar asociada es:
          \begin{equation}
              u_B = \frac{0{,}005 \times \bar{P}}{\sqrt{3}}
                  = \PLACEHOLDER{X{,}X}~\text{W}
              \label{eq:incertidumbre_B}
          \end{equation}

    \item \textbf{Tipo~B — variación de la tensión de red:}
          fluctuaciones de la tensión de alimentación durante el
          ensayo se propagan a la medición de potencia.
          Con una variación observada de $\Delta V = \PLACEHOLDER{X{,}X}$~V
          y el coeficiente de sensibilidad
          $c_V = \partial P / \partial V \approx I_{\text{rms}}$:
          \begin{equation}
              u_{V} = c_V \cdot \frac{\Delta V}{\sqrt{3}}
                    = \PLACEHOLDER{X{,}X}~\text{W}
              \label{eq:incertidumbre_V}
          \end{equation}

    \item \textbf{Tipo~B — resolución del sensor:}
          la resolución en potencia del PZEM-004T es 0{,}1~W.
          Asumiendo distribución rectangular:
          \begin{equation}
              u_{\text{res}} = \frac{0{,}1}{2\sqrt{3}}
                             \approx 0{,}029~\text{W}
              \label{eq:incertidumbre_res}
          \end{equation}
\end{itemize}

\subsubsection{Incertidumbre combinada y expandida}

Considerando las fuentes anteriores como independientes entre sí,
la incertidumbre estándar combinada es:

\begin{equation}
    u_c = \sqrt{u_{\bar{P}}^2 + u_B^2 + u_V^2 + u_{\text{res}}^2}
        = \PLACEHOLDER{X{,}X}~\text{W}
    \label{eq:incertidumbre_combinada}
\end{equation}

La incertidumbre expandida, con factor de cobertura $k=2$
(nivel de confianza $\approx$95\%, distribución normal), es:

\begin{equation}
    U = k \cdot u_c = 2 \times \PLACEHOLDER{X{,}X}
      = \PLACEHOLDER{X{,}X}~\text{W}
    \label{eq:incertidumbre_expandida}
\end{equation}

El resultado final de la medición de potencia se expresa como:

\begin{equation}
    P = (\PLACEHOLDER{XXXX{,}X} \pm \PLACEHOLDER{X{,}X})~\text{W}
        \quad (k=2,\; p \approx 95\%)
    \label{eq:resultado_final}
\end{equation}

\subsubsection{Relación entre incertidumbre y límite normativo}

La incertidumbre expandida del sistema ($U = \PLACEHOLDER{X{,}X}$~W)
representa el \PLACEHOLDER{X{,}X}\% del límite normativo
($\pm$10\% de $P_{\text{dec}} = \PLACEHOLDER{XXXX}$~W).
Este valor es \PLACEHOLDER{aceptable para los propósitos del
prototipo / recomendable reducir mediante calibración del sensor},
dado que el criterio general de metrología establece que la
incertidumbre del instrumento de medición debe ser al menos
cuatro veces menor que la tolerancia que se desea verificar.

% ------------------------------------------------------------
\subsection{Evaluación de la repetibilidad del sistema}
\label{subsec:repetibilidad}
% ------------------------------------------------------------

La repetibilidad del sistema se evalúa como la capacidad de
obtener resultados consistentes bajo las mismas condiciones
experimentales en mediciones sucesivas.

\begin{table}[H]
\centering
\caption{Indicadores de repetibilidad del ensayo de potencia}
\label{tab:repetibilidad}
\renewcommand{\arraystretch}{1.35}
\begin{tabular}{@{} l r l @{}}
\toprule
\textbf{Indicador} & \textbf{Valor} & \textbf{Interpretación} \\
\midrule
Media $\bar{P}$                 & \PLACEHOLDER{XXXX,X} W  & Valor central de la distribución \\
Desviación estándar $s_P$       & \PLACEHOLDER{X,X} W     & Dispersión de las mediciones \\
Coeficiente de variación CV     & \PLACEHOLDER{X,X}\%     & Repetibilidad relativa \\
Rango ($P_{\max} - P_{\min}$)   & \PLACEHOLDER{X,X} W     & Amplitud total de dispersión \\
Incertidumbre estándar $u_{\bar{P}}$ & \PLACEHOLDER{X,X} W & Incertidumbre tipo~A \\
\midrule
\textbf{Criterio de repetibilidad} & CV $<$ 2\%
    & \PLACEHOLDER{\textcolor{green!60!black}{\textbf{CUMPLE}} /
                   \textcolor{red}{\textbf{NO CUMPLE}}} \\
\bottomrule
\end{tabular}
\end{table}

% ------------------------------------------------------------
\subsection{Discusión de resultados}
\label{subsec:discusion}
% ------------------------------------------------------------

Los resultados obtenidos permiten extraer las siguientes
conclusiones sobre el desempeño del prototipo:

\paragraph{Exactitud en la medición de potencia.}
El error relativo promedio frente al instrumento de referencia
fue de \PLACEHOLDER{X{,}XX}\%, valor que
\PLACEHOLDER{se encuentra dentro / supera levemente}
la especificación del fabricante del módulo PZEM-004T (±0{,}5\%).
\PLACEHOLDER{Describir aquí la causa probable del error observado
si supera la especificación, p.ej.: caída de tensión en el cableado,
temperatura de operación del módulo, saturación del TC, etc.}

\paragraph{Repetibilidad.}
El coeficiente de variación obtenido de \PLACEHOLDER{X{,}X}\%
\PLACEHOLDER{confirma / no confirma} una repetibilidad adecuada
del sistema para el propósito del ensayo.
\PLACEHOLDER{Si CV > 2\%: identificar causa de la variabilidad
(fluctuaciones de red, termostato del DUT activo durante la medición,
interferencias en el bus UART, etc.) y proponer acción correctiva.}

\paragraph{Cumplimiento normativo.}
El DUT ensayado \PLACEHOLDER{aprobó / reprobó} el criterio de
consumo de la norma IEC~60335-1 (cláusula~10.2), con una desviación
de \PLACEHOLDER{X{,}X}\% respecto al valor declarado,
dentro del límite admisible de $\pm$10\%.
Este resultado \PLACEHOLDER{es / no es} consistente con lo esperado
para un electrodoméstico \PLACEHOLDER{en buenas condiciones / con
el defecto identificado: descripción}.

\paragraph{Limitaciones identificadas.}
Durante la campaña de ensayos se identificaron las siguientes
limitaciones del prototipo en su estado actual:

\begin{itemize}
    \item \PLACEHOLDER{Limitación 1: p.ej., el módulo PZEM-004T
          presenta lecturas inestables durante los primeros 30~s
          de conexión, por lo que se implementó un retardo de
          estabilización en el software.}
    \item \PLACEHOLDER{Limitación 2: p.ej., la precisión de la
          medición de factor de potencia se degrada para cargas
          muy resistivas ($\cos\varphi > 0{,}98$) debido a la
          resolución de 0{,}01 del módulo.}
    \item \PLACEHOLDER{Limitación 3: agregar según observaciones
          reales del ensayo.}
\end{itemize}

Estas limitaciones no comprometen la validez del ensayo de consumo
como función principal del prototipo, pero deben ser consideradas
al momento de ampliar el sistema hacia los ensayos de fuga y
rigidez dieléctrica contemplados en la arquitectura del equipo.

% ------------------------------------------------------------
\subsection{Comparación con el criterio de éxito del proyecto}
\label{subsec:criterio_exito}
% ------------------------------------------------------------

La Tabla~\ref{tab:criterios_exito} evalúa los resultados obtenidos
frente a los criterios de éxito definidos en la
Sección~\ref{sec:criterios_evaluacion} del presente informe.

\begin{table}[H]
\centering
\caption{Evaluación de criterios de éxito del prototipo}
\label{tab:criterios_exito}
\renewcommand{\arraystretch}{1.4}
\begin{tabular}{@{} p{5cm} p{3cm} p{3cm} p{1.8cm} @{}}
\toprule
\textbf{Criterio} & \textbf{Meta definida} & \textbf{Resultado obtenido} & \textbf{Estado} \\
\midrule
Reducción del tiempo de prueba vs.\ método manual
    & $\geq 30\%$
    & \PLACEHOLDER{XX\%}
    & \PLACEHOLDER{\textcolor{green!60!black}{\checkmark} /
                   \textcolor{red}{$\times$}} \\
\addlinespace
Reducción de intervención del operario
    & $\geq 80\%$
    & \PLACEHOLDER{XX\%}
    & \PLACEHOLDER{\textcolor{green!60!black}{\checkmark} /
                   \textcolor{red}{$\times$}} \\
\addlinespace
Error relativo en potencia frente a patrón
    & Según especif.\ PZEM ($\pm$0{,}5\%)
    & \PLACEHOLDER{$\pm$X{,}XX\%}
    & \PLACEHOLDER{\textcolor{green!60!black}{\checkmark} /
                   \textcolor{red}{$\times$}} \\
\addlinespace
Cumplimiento normativo (criterio IEC~10.2)
    & Detección correcta APROBADO / REPROBADO
    & \PLACEHOLDER{Correcto / Incorrecto}
    & \PLACEHOLDER{\textcolor{green!60!black}{\checkmark} /
                   \textcolor{red}{$\times$}} \\
\addlinespace
Repetibilidad (CV)
    & CV $< 2\%$
    & \PLACEHOLDER{X{,}X\%}
    & \PLACEHOLDER{\textcolor{green!60!black}{\checkmark} /
                   \textcolor{red}{$\times$}} \\
\addlinespace
Ajuste al presupuesto
    & $\leq$ USD~1030
    & USD~\PLACEHOLDER{XXX}
    & \PLACEHOLDER{\textcolor{green!60!black}{\checkmark} /
                   \textcolor{red}{$\times$}} \\
\bottomrule
\end{tabular}
\end{table}

% ============================================================

% Comando auxiliar para resaltar placeholders visualmente en Overleaf
% (comentar o eliminar antes de la entrega final)
\newcommand{\PLACEHOLDER}[1]{\textcolor{red}{\textbf{[#1]}}}

% ============================================================
\section{Resultados y Análisis}
\label{sec:resultados}
% ============================================================

Esta sección presenta los resultados obtenidos durante la campaña
de ensayos de validación del prototipo desarrollado.
Se exponen las mediciones realizadas sobre el dispositivo bajo
prueba (DUT), se comparan con los valores de referencia
establecidos por la norma IEC~60335-1 y con un instrumento patrón,
y se analiza el desempeño del sistema en términos de exactitud,
repetibilidad y cumplimiento normativo.

% ------------------------------------------------------------
\subsection{Configuración experimental}
\label{subsec:config_experimental}
% ------------------------------------------------------------

\subsubsection{Dispositivo bajo prueba (DUT)}

Las mediciones se realizaron sobre un electrodoméstico de uso
doméstico de la categoría \textbf{\PLACEHOLDER{nombre del aparato,
p.ej.: plancha a vapor}}, cuyas características declaradas por
el fabricante se resumen en la Tabla~\ref{tab:dut_datos}.

\begin{table}[H]
\centering
\caption{Datos de placa del dispositivo bajo prueba}
\label{tab:dut_datos}
\renewcommand{\arraystretch}{1.35}
\begin{tabular}{@{} l l @{}}
\toprule
\textbf{Parámetro} & \textbf{Valor declarado} \\
\midrule
Marca / modelo          & \PLACEHOLDER{Marca — Modelo} \\
Clase de aislamiento    & \PLACEHOLDER{Clase I / Clase II} \\
Tensión nominal         & \PLACEHOLDER{220 / 230} V AC, 50 Hz \\
Potencia nominal        & \PLACEHOLDER{XXX} W \\
Corriente nominal       & \PLACEHOLDER{X{,}X} A \\
Factor de potencia      & \PLACEHOLDER{X{,}XX} (si figura en placa) \\
Norma de producto       & \PLACEHOLDER{IEC 60335-2-3 (planchas) u otra} \\
\bottomrule
\end{tabular}
\end{table}

\subsubsection{Instrumentación utilizada}

El ensayo se ejecutó con el prototipo desarrollado en este proyecto.
Adicionalmente, como instrumento de referencia para la validación
de las mediciones, se empleó \PLACEHOLDER{nombre y modelo del
instrumento patrón, p.ej.: analizador de potencia Fluke 435-II /
multímetro Fluke 87V / vatímetro patrón del laboratorio}.

Las especificaciones de exactitud del instrumento patrón se
detallan en la Tabla~\ref{tab:instrumento_patron}.

\begin{table}[H]
\centering
\caption{Especificaciones del instrumento de referencia utilizado}
\label{tab:instrumento_patron}
\renewcommand{\arraystretch}{1.35}
\begin{tabular}{@{} l l @{}}
\toprule
\textbf{Parámetro} & \textbf{Especificación} \\
\midrule
Marca / modelo              & \PLACEHOLDER{Marca — Modelo} \\
Exactitud en tensión        & \PLACEHOLDER{$\pm$X{,}X\% + X dígitos} \\
Exactitud en corriente      & \PLACEHOLDER{$\pm$X{,}X\% + X dígitos} \\
Exactitud en potencia       & \PLACEHOLDER{$\pm$X{,}X\%} \\
Trazabilidad de calibración & \PLACEHOLDER{Calibrado por laboratorio
                              acreditado OAA N\textsuperscript{o}~XXX,
                              fecha: MM/AAAA} \\
\bottomrule
\end{tabular}
\end{table}

\subsubsection{Condiciones ambientales}

Las condiciones del laboratorio durante la realización de los
ensayos se registraron al inicio y al final de la sesión de
mediciones y se presentan en la Tabla~\ref{tab:condiciones_amb}.

\begin{table}[H]
\centering
\caption{Condiciones ambientales durante el ensayo}
\label{tab:condiciones_amb}
\renewcommand{\arraystretch}{1.35}
\begin{tabular}{@{} l r r r @{}}
\toprule
\textbf{Parámetro} & \textbf{Inicio} & \textbf{Final} & \textbf{Límite IEC} \\
\midrule
Temperatura ambiente [°C]   & \PLACEHOLDER{XX,X} & \PLACEHOLDER{XX,X} & $(20 \pm 5)$~°C \\
Humedad relativa [\%]        & \PLACEHOLDER{XX}   & \PLACEHOLDER{XX}   & $< 85$\% \\
Tensión de red [V$_{\text{rms}}$] & \PLACEHOLDER{XXX,X} & \PLACEHOLDER{XXX,X} & $V_n \pm 1\%$ \\
Frecuencia de red [Hz]       & \PLACEHOLDER{XX,X} & \PLACEHOLDER{XX,X} & $50 \pm 0{,}5$~Hz \\
\bottomrule
\end{tabular}
\end{table}

% ------------------------------------------------------------
\subsection{Resultados del ensayo de consumo de potencia}
\label{subsec:resultados_potencia}
% ------------------------------------------------------------

\subsubsection{Mediciones individuales}

Se realizaron \PLACEHOLDER{N} mediciones consecutivas
con el DUT en condición térmica estable
(temperatura de operación normal alcanzada tras
\PLACEHOLDER{X}~minutos de precalentamiento).
Cada medición corresponde a un ciclo completo de adquisición
del módulo PZEM-004T~v3.0, con un intervalo de muestreo de
\PLACEHOLDER{X}~segundos.
Los valores registrados se presentan en la
Tabla~\ref{tab:mediciones_pzem}.

\begin{table}[H]
\centering
\caption{Mediciones de parámetros eléctricos — DUT: \PLACEHOLDER{nombre del aparato}}
\label{tab:mediciones_pzem}
\renewcommand{\arraystretch}{1.3}
\rowcolors{2}{gray!8}{white}
\begin{tabular}{@{} c r r r r r @{}}
\toprule
\textbf{Muestra} &
\textbf{$V_{\text{rms}}$ [V]} &
\textbf{$I_{\text{rms}}$ [A]} &
\textbf{$P$ [W]} &
\textbf{$\cos\varphi$} &
\textbf{$f$ [Hz]} \\
\midrule
1  & \PLACEHOLDER{XXX,X} & \PLACEHOLDER{X,XXX} & \PLACEHOLDER{XXXX,X} & \PLACEHOLDER{X,XX} & \PLACEHOLDER{XX,X} \\
2  & \PLACEHOLDER{XXX,X} & \PLACEHOLDER{X,XXX} & \PLACEHOLDER{XXXX,X} & \PLACEHOLDER{X,XX} & \PLACEHOLDER{XX,X} \\
3  & \PLACEHOLDER{XXX,X} & \PLACEHOLDER{X,XXX} & \PLACEHOLDER{XXXX,X} & \PLACEHOLDER{X,XX} & \PLACEHOLDER{XX,X} \\
4  & \PLACEHOLDER{XXX,X} & \PLACEHOLDER{X,XXX} & \PLACEHOLDER{XXXX,X} & \PLACEHOLDER{X,XX} & \PLACEHOLDER{XX,X} \\
5  & \PLACEHOLDER{XXX,X} & \PLACEHOLDER{X,XXX} & \PLACEHOLDER{XXXX,X} & \PLACEHOLDER{X,XX} & \PLACEHOLDER{XX,X} \\
6  & \PLACEHOLDER{XXX,X} & \PLACEHOLDER{X,XXX} & \PLACEHOLDER{XXXX,X} & \PLACEHOLDER{X,XX} & \PLACEHOLDER{XX,X} \\
7  & \PLACEHOLDER{XXX,X} & \PLACEHOLDER{X,XXX} & \PLACEHOLDER{XXXX,X} & \PLACEHOLDER{X,XX} & \PLACEHOLDER{XX,X} \\
8  & \PLACEHOLDER{XXX,X} & \PLACEHOLDER{X,XXX} & \PLACEHOLDER{XXXX,X} & \PLACEHOLDER{X,XX} & \PLACEHOLDER{XX,X} \\
9  & \PLACEHOLDER{XXX,X} & \PLACEHOLDER{X,XXX} & \PLACEHOLDER{XXXX,X} & \PLACEHOLDER{X,XX} & \PLACEHOLDER{XX,X} \\
10 & \PLACEHOLDER{XXX,X} & \PLACEHOLDER{X,XXX} & \PLACEHOLDER{XXXX,X} & \PLACEHOLDER{X,XX} & \PLACEHOLDER{XX,X} \\
\midrule
\textbf{Media} $\bar{x}$
   & \PLACEHOLDER{XXX,X} & \PLACEHOLDER{X,XXX} & \PLACEHOLDER{XXXX,X} & \PLACEHOLDER{X,XX} & \PLACEHOLDER{XX,X} \\
\textbf{Desv. est.} $s$
   & \PLACEHOLDER{X,X}   & \PLACEHOLDER{X,XXX} & \PLACEHOLDER{X,X}    & \PLACEHOLDER{X,XX} & \PLACEHOLDER{X,X}  \\
\textbf{Mín.}
   & \PLACEHOLDER{XXX,X} & \PLACEHOLDER{X,XXX} & \PLACEHOLDER{XXXX,X} & \PLACEHOLDER{X,XX} & \PLACEHOLDER{XX,X} \\
\textbf{Máx.}
   & \PLACEHOLDER{XXX,X} & \PLACEHOLDER{X,XXX} & \PLACEHOLDER{XXXX,X} & \PLACEHOLDER{X,XX} & \PLACEHOLDER{XX,X} \\
\bottomrule
\end{tabular}
\end{table}

\subsubsection{Análisis estadístico}

A partir de las \PLACEHOLDER{N} mediciones registradas,
se calcularon los estadísticos descriptivos de la potencia activa $P$:

\begin{align}
    \bar{P} &= \frac{1}{N}\sum_{i=1}^{N} P_i
             = \PLACEHOLDER{XXXX{,}X}~\text{W}
    \label{eq:media_potencia} \\[6pt]
    s_P &= \sqrt{\frac{1}{N-1}\sum_{i=1}^{N}(P_i - \bar{P})^2}
          = \PLACEHOLDER{X{,}X}~\text{W}
    \label{eq:std_potencia} \\[6pt]
    u_{\bar{P}} &= \frac{s_P}{\sqrt{N}}
                 = \PLACEHOLDER{X{,}X}~\text{W}
    \label{eq:incertidumbre_estadistica}
\end{align}

\noindent donde $u_{\bar{P}}$ es la incertidumbre estándar tipo~A
asociada a la dispersión de las mediciones.

El \textbf{coeficiente de variación} (CV) de la potencia medida es:

\begin{equation}
    \text{CV} = \frac{s_P}{\bar{P}} \times 100\%
              = \PLACEHOLDER{X{,}X}\%
    \label{eq:cv_potencia}
\end{equation}

Un CV menor al 2\% indica una buena repetibilidad del sistema
de medición bajo las condiciones de ensayo.

% ------------------------------------------------------------
\subsection{Validación frente al instrumento de referencia}
\label{subsec:validacion_patron}
% ------------------------------------------------------------

\subsubsection{Comparación de mediciones}

Con el objetivo de verificar la exactitud del prototipo,
se realizaron mediciones simultáneas con el equipo desarrollado
y con el instrumento de referencia sobre el mismo DUT y en las
mismas condiciones de carga.
Los resultados comparativos se presentan en la
Tabla~\ref{tab:comparacion_patron}.

\begin{table}[H]
\centering
\caption{Comparación entre el prototipo y el instrumento de referencia}
\label{tab:comparacion_patron}
\renewcommand{\arraystretch}{1.4}
\begin{tabular}{@{} l r r r r @{}}
\toprule
\textbf{Parámetro}
    & \textbf{Prototipo}
    & \textbf{Referencia}
    & \textbf{Error abs.}
    & \textbf{Error rel. [\%]} \\
\midrule
$V_{\text{rms}}$ [V]
    & \PLACEHOLDER{XXX,X}
    & \PLACEHOLDER{XXX,X}
    & \PLACEHOLDER{$\pm$X,X}
    & \PLACEHOLDER{X,XX} \\
$I_{\text{rms}}$ [A]
    & \PLACEHOLDER{X,XXX}
    & \PLACEHOLDER{X,XXX}
    & \PLACEHOLDER{$\pm$X,XXX}
    & \PLACEHOLDER{X,XX} \\
$P$ [W]
    & \PLACEHOLDER{XXXX,X}
    & \PLACEHOLDER{XXXX,X}
    & \PLACEHOLDER{$\pm$X,X}
    & \PLACEHOLDER{X,XX} \\
$\cos\varphi$
    & \PLACEHOLDER{X,XX}
    & \PLACEHOLDER{X,XX}
    & \PLACEHOLDER{$\pm$X,XX}
    & \PLACEHOLDER{X,XX} \\
\bottomrule
\end{tabular}
\end{table}

\noindent El error relativo en potencia activa $P$ se calcula como:

\begin{equation}
    e_r = \frac{P_{\text{prototipo}} - P_{\text{referencia}}}
               {P_{\text{referencia}}} \times 100\%
        = \PLACEHOLDER{X{,}XX}\%
    \label{eq:error_relativo}
\end{equation}

La especificación de exactitud del módulo PZEM-004T~v3.0 en
potencia activa es de $\pm$0{,}5\%.
El error relativo obtenido de \PLACEHOLDER{X{,}XX}\%
\PLACEHOLDER{es / no es} consistente con dicha especificación,
\PLACEHOLDER{lo que confirma / lo que sugiere una posible causa
de error en: descripción del problema observado}.

\subsubsection{Cumplimiento del criterio normativo}

La norma IEC~60335-1 (cláusula~10.2) establece que la potencia
consumida por el DUT no debe diferir de la potencia declarada
por el fabricante en más de $\pm$10\%.
La evaluación de este criterio se resume en la
Tabla~\ref{tab:criterio_normativo}.

\begin{table}[H]
\centering
\caption{Evaluación del criterio normativo de consumo — IEC~60335-1, cláusula~10.2}
\label{tab:criterio_normativo}
\renewcommand{\arraystretch}{1.4}
\begin{tabular}{@{} l r @{}}
\toprule
\textbf{Parámetro} & \textbf{Valor} \\
\midrule
Potencia declarada por el fabricante ($P_{\text{dec}}$)
    & \PLACEHOLDER{XXXX} W \\
Límite inferior admisible ($P_{\text{dec}} \times 0{,}90$)
    & \PLACEHOLDER{XXXX} W \\
Límite superior admisible ($P_{\text{dec}} \times 1{,}10$)
    & \PLACEHOLDER{XXXX} W \\
Potencia media medida por el prototipo ($\bar{P}$)
    & \PLACEHOLDER{XXXX,X} W \\
Desviación respecto al valor declarado
    & \PLACEHOLDER{$+$X{,}X / $-$X{,}X}\% \\
\midrule
\textbf{Resultado del ensayo}
    & \PLACEHOLDER{\textcolor{green!60!black}{\textbf{APROBADO}} /
                   \textcolor{red}{\textbf{REPROBADO}}} \\
\bottomrule
\end{tabular}
\end{table}

% ------------------------------------------------------------
\subsection{Análisis de incertidumbre de medición}
\label{subsec:incertidumbre}
% ------------------------------------------------------------

El análisis de incertidumbre se realiza siguiendo el marco de la
Guía para la Expresión de la Incertidumbre de Medición
(GUM, JCGM~100:2008), considerando las fuentes de incertidumbre
relevantes para la medición de potencia activa.

\subsubsection{Fuentes de incertidumbre identificadas}

\begin{itemize}
    \item \textbf{Tipo~A — dispersión estadística de las mediciones:}
          incertidumbre debida a la variabilidad aleatoria del proceso
          de medición, cuantificada por la desviación estándar de la
          media $u_{\bar{P}}$ (ecuación~\ref{eq:incertidumbre_estadistica}).

    \item \textbf{Tipo~B — exactitud del sensor PZEM-004T:}
          la hoja de datos del módulo especifica una exactitud de
          $\pm$0{,}5\% en potencia activa.
          Asumiendo distribución rectangular, la incertidumbre
          estándar asociada es:
          \begin{equation}
              u_B = \frac{0{,}005 \times \bar{P}}{\sqrt{3}}
                  = \PLACEHOLDER{X{,}X}~\text{W}
              \label{eq:incertidumbre_B}
          \end{equation}

    \item \textbf{Tipo~B — variación de la tensión de red:}
          fluctuaciones de la tensión de alimentación durante el
          ensayo se propagan a la medición de potencia.
          Con una variación observada de $\Delta V = \PLACEHOLDER{X{,}X}$~V
          y el coeficiente de sensibilidad
          $c_V = \partial P / \partial V \approx I_{\text{rms}}$:
          \begin{equation}
              u_{V} = c_V \cdot \frac{\Delta V}{\sqrt{3}}
                    = \PLACEHOLDER{X{,}X}~\text{W}
              \label{eq:incertidumbre_V}
          \end{equation}

    \item \textbf{Tipo~B — resolución del sensor:}
          la resolución en potencia del PZEM-004T es 0{,}1~W.
          Asumiendo distribución rectangular:
          \begin{equation}
              u_{\text{res}} = \frac{0{,}1}{2\sqrt{3}}
                             \approx 0{,}029~\text{W}
              \label{eq:incertidumbre_res}
          \end{equation}
\end{itemize}

\subsubsection{Incertidumbre combinada y expandida}

Considerando las fuentes anteriores como independientes entre sí,
la incertidumbre estándar combinada es:

\begin{equation}
    u_c = \sqrt{u_{\bar{P}}^2 + u_B^2 + u_V^2 + u_{\text{res}}^2}
        = \PLACEHOLDER{X{,}X}~\text{W}
    \label{eq:incertidumbre_combinada}
\end{equation}

La incertidumbre expandida, con factor de cobertura $k=2$
(nivel de confianza $\approx$95\%, distribución normal), es:

\begin{equation}
    U = k \cdot u_c = 2 \times \PLACEHOLDER{X{,}X}
      = \PLACEHOLDER{X{,}X}~\text{W}
    \label{eq:incertidumbre_expandida}
\end{equation}

El resultado final de la medición de potencia se expresa como:

\begin{equation}
    P = (\PLACEHOLDER{XXXX{,}X} \pm \PLACEHOLDER{X{,}X})~\text{W}
        \quad (k=2,\; p \approx 95\%)
    \label{eq:resultado_final}
\end{equation}

\subsubsection{Relación entre incertidumbre y límite normativo}

La incertidumbre expandida del sistema ($U = \PLACEHOLDER{X{,}X}$~W)
representa el \PLACEHOLDER{X{,}X}\% del límite normativo
($\pm$10\% de $P_{\text{dec}} = \PLACEHOLDER{XXXX}$~W).
Este valor es \PLACEHOLDER{aceptable para los propósitos del
prototipo / recomendable reducir mediante calibración del sensor},
dado que el criterio general de metrología establece que la
incertidumbre del instrumento de medición debe ser al menos
cuatro veces menor que la tolerancia que se desea verificar.

% ------------------------------------------------------------
\subsection{Evaluación de la repetibilidad del sistema}
\label{subsec:repetibilidad}
% ------------------------------------------------------------

La repetibilidad del sistema se evalúa como la capacidad de
obtener resultados consistentes bajo las mismas condiciones
experimentales en mediciones sucesivas.

\begin{table}[H]
\centering
\caption{Indicadores de repetibilidad del ensayo de potencia}
\label{tab:repetibilidad}
\renewcommand{\arraystretch}{1.35}
\begin{tabular}{@{} l r l @{}}
\toprule
\textbf{Indicador} & \textbf{Valor} & \textbf{Interpretación} \\
\midrule
Media $\bar{P}$                 & \PLACEHOLDER{XXXX,X} W  & Valor central de la distribución \\
Desviación estándar $s_P$       & \PLACEHOLDER{X,X} W     & Dispersión de las mediciones \\
Coeficiente de variación CV     & \PLACEHOLDER{X,X}\%     & Repetibilidad relativa \\
Rango ($P_{\max} - P_{\min}$)   & \PLACEHOLDER{X,X} W     & Amplitud total de dispersión \\
Incertidumbre estándar $u_{\bar{P}}$ & \PLACEHOLDER{X,X} W & Incertidumbre tipo~A \\
\midrule
\textbf{Criterio de repetibilidad} & CV $<$ 2\%
    & \PLACEHOLDER{\textcolor{green!60!black}{\textbf{CUMPLE}} /
                   \textcolor{red}{\textbf{NO CUMPLE}}} \\
\bottomrule
\end{tabular}
\end{table}

% ------------------------------------------------------------
\subsection{Discusión de resultados}
\label{subsec:discusion}
% ------------------------------------------------------------

Los resultados obtenidos permiten extraer las siguientes
conclusiones sobre el desempeño del prototipo:

\paragraph{Exactitud en la medición de potencia.}
El error relativo promedio frente al instrumento de referencia
fue de \PLACEHOLDER{X{,}XX}\%, valor que
\PLACEHOLDER{se encuentra dentro / supera levemente}
la especificación del fabricante del módulo PZEM-004T (±0{,}5\%).
\PLACEHOLDER{Describir aquí la causa probable del error observado
si supera la especificación, p.ej.: caída de tensión en el cableado,
temperatura de operación del módulo, saturación del TC, etc.}

\paragraph{Repetibilidad.}
El coeficiente de variación obtenido de \PLACEHOLDER{X{,}X}\%
\PLACEHOLDER{confirma / no confirma} una repetibilidad adecuada
del sistema para el propósito del ensayo.
\PLACEHOLDER{Si CV > 2\%: identificar causa de la variabilidad
(fluctuaciones de red, termostato del DUT activo durante la medición,
interferencias en el bus UART, etc.) y proponer acción correctiva.}

\paragraph{Cumplimiento normativo.}
El DUT ensayado \PLACEHOLDER{aprobó / reprobó} el criterio de
consumo de la norma IEC~60335-1 (cláusula~10.2), con una desviación
de \PLACEHOLDER{X{,}X}\% respecto al valor declarado,
dentro del límite admisible de $\pm$10\%.
Este resultado \PLACEHOLDER{es / no es} consistente con lo esperado
para un electrodoméstico \PLACEHOLDER{en buenas condiciones / con
el defecto identificado: descripción}.

\paragraph{Limitaciones identificadas.}
Durante la campaña de ensayos se identificaron las siguientes
limitaciones del prototipo en su estado actual:

\begin{itemize}
    \item \PLACEHOLDER{Limitación 1: p.ej., el módulo PZEM-004T
          presenta lecturas inestables durante los primeros 30~s
          de conexión, por lo que se implementó un retardo de
          estabilización en el software.}
    \item \PLACEHOLDER{Limitación 2: p.ej., la precisión de la
          medición de factor de potencia se degrada para cargas
          muy resistivas ($\cos\varphi > 0{,}98$) debido a la
          resolución de 0{,}01 del módulo.}
    \item \PLACEHOLDER{Limitación 3: agregar según observaciones
          reales del ensayo.}
\end{itemize}

Estas limitaciones no comprometen la validez del ensayo de consumo
como función principal del prototipo, pero deben ser consideradas
al momento de ampliar el sistema hacia los ensayos de fuga y
rigidez dieléctrica contemplados en la arquitectura del equipo.

% ------------------------------------------------------------
\subsection{Comparación con el criterio de éxito del proyecto}
\label{subsec:criterio_exito}
% ------------------------------------------------------------

La Tabla~\ref{tab:criterios_exito} evalúa los resultados obtenidos
frente a los criterios de éxito definidos en la
Sección~\ref{sec:criterios_evaluacion} del presente informe.

\begin{table}[H]
\centering
\caption{Evaluación de criterios de éxito del prototipo}
\label{tab:criterios_exito}
\renewcommand{\arraystretch}{1.4}
\begin{tabular}{@{} p{5cm} p{3cm} p{3cm} p{1.8cm} @{}}
\toprule
\textbf{Criterio} & \textbf{Meta definida} & \textbf{Resultado obtenido} & \textbf{Estado} \\
\midrule
Reducción del tiempo de prueba vs.\ método manual
    & $\geq 30\%$
    & \PLACEHOLDER{XX\%}
    & \PLACEHOLDER{\textcolor{green!60!black}{\checkmark} /
                   \textcolor{red}{$\times$}} \\
\addlinespace
Reducción de intervención del operario
    & $\geq 80\%$
    & \PLACEHOLDER{XX\%}
    & \PLACEHOLDER{\textcolor{green!60!black}{\checkmark} /
                   \textcolor{red}{$\times$}} \\
\addlinespace
Error relativo en potencia frente a patrón
    & Según especif.\ PZEM ($\pm$0{,}5\%)
    & \PLACEHOLDER{$\pm$X{,}XX\%}
    & \PLACEHOLDER{\textcolor{green!60!black}{\checkmark} /
                   \textcolor{red}{$\times$}} \\
\addlinespace
Cumplimiento normativo (criterio IEC~10.2)
    & Detección correcta APROBADO / REPROBADO
    & \PLACEHOLDER{Correcto / Incorrecto}
    & \PLACEHOLDER{\textcolor{green!60!black}{\checkmark} /
                   \textcolor{red}{$\times$}} \\
\addlinespace
Repetibilidad (CV)
    & CV $< 2\%$
    & \PLACEHOLDER{X{,}X\%}
    & \PLACEHOLDER{\textcolor{green!60!black}{\checkmark} /
                   \textcolor{red}{$\times$}} \\
\addlinespace
Ajuste al presupuesto
    & $\leq$ USD~1030
    & USD~\PLACEHOLDER{XXX}
    & \PLACEHOLDER{\textcolor{green!60!black}{\checkmark} /
                   \textcolor{red}{$\times$}} \\
\bottomrule
\end{tabular}
\end{table}

% ============================================================

% Comando auxiliar para resaltar placeholders visualmente en Overleaf
% (comentar o eliminar antes de la entrega final)
\newcommand{\PLACEHOLDER}[1]{\textcolor{red}{\textbf{[#1]}}}

% ============================================================
\section{Resultados y Análisis}
\label{sec:resultados}
% ============================================================

Esta sección presenta los resultados obtenidos durante la campaña
de ensayos de validación del prototipo desarrollado.
Se exponen las mediciones realizadas sobre el dispositivo bajo
prueba (DUT), se comparan con los valores de referencia
establecidos por la norma IEC~60335-1 y con un instrumento patrón,
y se analiza el desempeño del sistema en términos de exactitud,
repetibilidad y cumplimiento normativo.

% ------------------------------------------------------------
\subsection{Configuración experimental}
\label{subsec:config_experimental}
% ------------------------------------------------------------

\subsubsection{Dispositivo bajo prueba (DUT)}

Las mediciones se realizaron sobre un electrodoméstico de uso
doméstico de la categoría \textbf{\PLACEHOLDER{nombre del aparato,
p.ej.: plancha a vapor}}, cuyas características declaradas por
el fabricante se resumen en la Tabla~\ref{tab:dut_datos}.

\begin{table}[H]
\centering
\caption{Datos de placa del dispositivo bajo prueba}
\label{tab:dut_datos}
\renewcommand{\arraystretch}{1.35}
\begin{tabular}{@{} l l @{}}
\toprule
\textbf{Parámetro} & \textbf{Valor declarado} \\
\midrule
Marca / modelo          & \PLACEHOLDER{Marca — Modelo} \\
Clase de aislamiento    & \PLACEHOLDER{Clase I / Clase II} \\
Tensión nominal         & \PLACEHOLDER{220 / 230} V AC, 50 Hz \\
Potencia nominal        & \PLACEHOLDER{XXX} W \\
Corriente nominal       & \PLACEHOLDER{X{,}X} A \\
Factor de potencia      & \PLACEHOLDER{X{,}XX} (si figura en placa) \\
Norma de producto       & \PLACEHOLDER{IEC 60335-2-3 (planchas) u otra} \\
\bottomrule
\end{tabular}
\end{table}

\subsubsection{Instrumentación utilizada}

El ensayo se ejecutó con el prototipo desarrollado en este proyecto.
Adicionalmente, como instrumento de referencia para la validación
de las mediciones, se empleó \PLACEHOLDER{nombre y modelo del
instrumento patrón, p.ej.: analizador de potencia Fluke 435-II /
multímetro Fluke 87V / vatímetro patrón del laboratorio}.

Las especificaciones de exactitud del instrumento patrón se
detallan en la Tabla~\ref{tab:instrumento_patron}.

\begin{table}[H]
\centering
\caption{Especificaciones del instrumento de referencia utilizado}
\label{tab:instrumento_patron}
\renewcommand{\arraystretch}{1.35}
\begin{tabular}{@{} l l @{}}
\toprule
\textbf{Parámetro} & \textbf{Especificación} \\
\midrule
Marca / modelo              & \PLACEHOLDER{Marca — Modelo} \\
Exactitud en tensión        & \PLACEHOLDER{$\pm$X{,}X\% + X dígitos} \\
Exactitud en corriente      & \PLACEHOLDER{$\pm$X{,}X\% + X dígitos} \\
Exactitud en potencia       & \PLACEHOLDER{$\pm$X{,}X\%} \\
Trazabilidad de calibración & \PLACEHOLDER{Calibrado por laboratorio
                              acreditado OAA N\textsuperscript{o}~XXX,
                              fecha: MM/AAAA} \\
\bottomrule
\end{tabular}
\end{table}

\subsubsection{Condiciones ambientales}

Las condiciones del laboratorio durante la realización de los
ensayos se registraron al inicio y al final de la sesión de
mediciones y se presentan en la Tabla~\ref{tab:condiciones_amb}.

\begin{table}[H]
\centering
\caption{Condiciones ambientales durante el ensayo}
\label{tab:condiciones_amb}
\renewcommand{\arraystretch}{1.35}
\begin{tabular}{@{} l r r r @{}}
\toprule
\textbf{Parámetro} & \textbf{Inicio} & \textbf{Final} & \textbf{Límite IEC} \\
\midrule
Temperatura ambiente [°C]   & \PLACEHOLDER{XX,X} & \PLACEHOLDER{XX,X} & $(20 \pm 5)$~°C \\
Humedad relativa [\%]        & \PLACEHOLDER{XX}   & \PLACEHOLDER{XX}   & $< 85$\% \\
Tensión de red [V$_{\text{rms}}$] & \PLACEHOLDER{XXX,X} & \PLACEHOLDER{XXX,X} & $V_n \pm 1\%$ \\
Frecuencia de red [Hz]       & \PLACEHOLDER{XX,X} & \PLACEHOLDER{XX,X} & $50 \pm 0{,}5$~Hz \\
\bottomrule
\end{tabular}
\end{table}

% ------------------------------------------------------------
\subsection{Resultados del ensayo de consumo de potencia}
\label{subsec:resultados_potencia}
% ------------------------------------------------------------

\subsubsection{Mediciones individuales}

Se realizaron \PLACEHOLDER{N} mediciones consecutivas
con el DUT en condición térmica estable
(temperatura de operación normal alcanzada tras
\PLACEHOLDER{X}~minutos de precalentamiento).
Cada medición corresponde a un ciclo completo de adquisición
del módulo PZEM-004T~v3.0, con un intervalo de muestreo de
\PLACEHOLDER{X}~segundos.
Los valores registrados se presentan en la
Tabla~\ref{tab:mediciones_pzem}.

\begin{table}[H]
\centering
\caption{Mediciones de parámetros eléctricos — DUT: \PLACEHOLDER{nombre del aparato}}
\label{tab:mediciones_pzem}
\renewcommand{\arraystretch}{1.3}
\rowcolors{2}{gray!8}{white}
\begin{tabular}{@{} c r r r r r @{}}
\toprule
\textbf{Muestra} &
\textbf{$V_{\text{rms}}$ [V]} &
\textbf{$I_{\text{rms}}$ [A]} &
\textbf{$P$ [W]} &
\textbf{$\cos\varphi$} &
\textbf{$f$ [Hz]} \\
\midrule
1  & \PLACEHOLDER{XXX,X} & \PLACEHOLDER{X,XXX} & \PLACEHOLDER{XXXX,X} & \PLACEHOLDER{X,XX} & \PLACEHOLDER{XX,X} \\
2  & \PLACEHOLDER{XXX,X} & \PLACEHOLDER{X,XXX} & \PLACEHOLDER{XXXX,X} & \PLACEHOLDER{X,XX} & \PLACEHOLDER{XX,X} \\
3  & \PLACEHOLDER{XXX,X} & \PLACEHOLDER{X,XXX} & \PLACEHOLDER{XXXX,X} & \PLACEHOLDER{X,XX} & \PLACEHOLDER{XX,X} \\
4  & \PLACEHOLDER{XXX,X} & \PLACEHOLDER{X,XXX} & \PLACEHOLDER{XXXX,X} & \PLACEHOLDER{X,XX} & \PLACEHOLDER{XX,X} \\
5  & \PLACEHOLDER{XXX,X} & \PLACEHOLDER{X,XXX} & \PLACEHOLDER{XXXX,X} & \PLACEHOLDER{X,XX} & \PLACEHOLDER{XX,X} \\
6  & \PLACEHOLDER{XXX,X} & \PLACEHOLDER{X,XXX} & \PLACEHOLDER{XXXX,X} & \PLACEHOLDER{X,XX} & \PLACEHOLDER{XX,X} \\
7  & \PLACEHOLDER{XXX,X} & \PLACEHOLDER{X,XXX} & \PLACEHOLDER{XXXX,X} & \PLACEHOLDER{X,XX} & \PLACEHOLDER{XX,X} \\
8  & \PLACEHOLDER{XXX,X} & \PLACEHOLDER{X,XXX} & \PLACEHOLDER{XXXX,X} & \PLACEHOLDER{X,XX} & \PLACEHOLDER{XX,X} \\
9  & \PLACEHOLDER{XXX,X} & \PLACEHOLDER{X,XXX} & \PLACEHOLDER{XXXX,X} & \PLACEHOLDER{X,XX} & \PLACEHOLDER{XX,X} \\
10 & \PLACEHOLDER{XXX,X} & \PLACEHOLDER{X,XXX} & \PLACEHOLDER{XXXX,X} & \PLACEHOLDER{X,XX} & \PLACEHOLDER{XX,X} \\
\midrule
\textbf{Media} $\bar{x}$
   & \PLACEHOLDER{XXX,X} & \PLACEHOLDER{X,XXX} & \PLACEHOLDER{XXXX,X} & \PLACEHOLDER{X,XX} & \PLACEHOLDER{XX,X} \\
\textbf{Desv. est.} $s$
   & \PLACEHOLDER{X,X}   & \PLACEHOLDER{X,XXX} & \PLACEHOLDER{X,X}    & \PLACEHOLDER{X,XX} & \PLACEHOLDER{X,X}  \\
\textbf{Mín.}
   & \PLACEHOLDER{XXX,X} & \PLACEHOLDER{X,XXX} & \PLACEHOLDER{XXXX,X} & \PLACEHOLDER{X,XX} & \PLACEHOLDER{XX,X} \\
\textbf{Máx.}
   & \PLACEHOLDER{XXX,X} & \PLACEHOLDER{X,XXX} & \PLACEHOLDER{XXXX,X} & \PLACEHOLDER{X,XX} & \PLACEHOLDER{XX,X} \\
\bottomrule
\end{tabular}
\end{table}

\subsubsection{Análisis estadístico}

A partir de las \PLACEHOLDER{N} mediciones registradas,
se calcularon los estadísticos descriptivos de la potencia activa $P$:

\begin{align}
    \bar{P} &= \frac{1}{N}\sum_{i=1}^{N} P_i
             = \PLACEHOLDER{XXXX{,}X}~\text{W}
    \label{eq:media_potencia} \\[6pt]
    s_P &= \sqrt{\frac{1}{N-1}\sum_{i=1}^{N}(P_i - \bar{P})^2}
          = \PLACEHOLDER{X{,}X}~\text{W}
    \label{eq:std_potencia} \\[6pt]
    u_{\bar{P}} &= \frac{s_P}{\sqrt{N}}
                 = \PLACEHOLDER{X{,}X}~\text{W}
    \label{eq:incertidumbre_estadistica}
\end{align}

\noindent donde $u_{\bar{P}}$ es la incertidumbre estándar tipo~A
asociada a la dispersión de las mediciones.

El \textbf{coeficiente de variación} (CV) de la potencia medida es:

\begin{equation}
    \text{CV} = \frac{s_P}{\bar{P}} \times 100\%
              = \PLACEHOLDER{X{,}X}\%
    \label{eq:cv_potencia}
\end{equation}

Un CV menor al 2\% indica una buena repetibilidad del sistema
de medición bajo las condiciones de ensayo.

% ------------------------------------------------------------
\subsection{Validación frente al instrumento de referencia}
\label{subsec:validacion_patron}
% ------------------------------------------------------------

\subsubsection{Comparación de mediciones}

Con el objetivo de verificar la exactitud del prototipo,
se realizaron mediciones simultáneas con el equipo desarrollado
y con el instrumento de referencia sobre el mismo DUT y en las
mismas condiciones de carga.
Los resultados comparativos se presentan en la
Tabla~\ref{tab:comparacion_patron}.

\begin{table}[H]
\centering
\caption{Comparación entre el prototipo y el instrumento de referencia}
\label{tab:comparacion_patron}
\renewcommand{\arraystretch}{1.4}
\begin{tabular}{@{} l r r r r @{}}
\toprule
\textbf{Parámetro}
    & \textbf{Prototipo}
    & \textbf{Referencia}
    & \textbf{Error abs.}
    & \textbf{Error rel. [\%]} \\
\midrule
$V_{\text{rms}}$ [V]
    & \PLACEHOLDER{XXX,X}
    & \PLACEHOLDER{XXX,X}
    & \PLACEHOLDER{$\pm$X,X}
    & \PLACEHOLDER{X,XX} \\
$I_{\text{rms}}$ [A]
    & \PLACEHOLDER{X,XXX}
    & \PLACEHOLDER{X,XXX}
    & \PLACEHOLDER{$\pm$X,XXX}
    & \PLACEHOLDER{X,XX} \\
$P$ [W]
    & \PLACEHOLDER{XXXX,X}
    & \PLACEHOLDER{XXXX,X}
    & \PLACEHOLDER{$\pm$X,X}
    & \PLACEHOLDER{X,XX} \\
$\cos\varphi$
    & \PLACEHOLDER{X,XX}
    & \PLACEHOLDER{X,XX}
    & \PLACEHOLDER{$\pm$X,XX}
    & \PLACEHOLDER{X,XX} \\
\bottomrule
\end{tabular}
\end{table}

\noindent El error relativo en potencia activa $P$ se calcula como:

\begin{equation}
    e_r = \frac{P_{\text{prototipo}} - P_{\text{referencia}}}
               {P_{\text{referencia}}} \times 100\%
        = \PLACEHOLDER{X{,}XX}\%
    \label{eq:error_relativo}
\end{equation}

La especificación de exactitud del módulo PZEM-004T~v3.0 en
potencia activa es de $\pm$0{,}5\%.
El error relativo obtenido de \PLACEHOLDER{X{,}XX}\%
\PLACEHOLDER{es / no es} consistente con dicha especificación,
\PLACEHOLDER{lo que confirma / lo que sugiere una posible causa
de error en: descripción del problema observado}.

\subsubsection{Cumplimiento del criterio normativo}

La norma IEC~60335-1 (cláusula~10.2) establece que la potencia
consumida por el DUT no debe diferir de la potencia declarada
por el fabricante en más de $\pm$10\%.
La evaluación de este criterio se resume en la
Tabla~\ref{tab:criterio_normativo}.

\begin{table}[H]
\centering
\caption{Evaluación del criterio normativo de consumo — IEC~60335-1, cláusula~10.2}
\label{tab:criterio_normativo}
\renewcommand{\arraystretch}{1.4}
\begin{tabular}{@{} l r @{}}
\toprule
\textbf{Parámetro} & \textbf{Valor} \\
\midrule
Potencia declarada por el fabricante ($P_{\text{dec}}$)
    & \PLACEHOLDER{XXXX} W \\
Límite inferior admisible ($P_{\text{dec}} \times 0{,}90$)
    & \PLACEHOLDER{XXXX} W \\
Límite superior admisible ($P_{\text{dec}} \times 1{,}10$)
    & \PLACEHOLDER{XXXX} W \\
Potencia media medida por el prototipo ($\bar{P}$)
    & \PLACEHOLDER{XXXX,X} W \\
Desviación respecto al valor declarado
    & \PLACEHOLDER{$+$X{,}X / $-$X{,}X}\% \\
\midrule
\textbf{Resultado del ensayo}
    & \PLACEHOLDER{\textcolor{green!60!black}{\textbf{APROBADO}} /
                   \textcolor{red}{\textbf{REPROBADO}}} \\
\bottomrule
\end{tabular}
\end{table}

% ------------------------------------------------------------
\subsection{Análisis de incertidumbre de medición}
\label{subsec:incertidumbre}
% ------------------------------------------------------------

El análisis de incertidumbre se realiza siguiendo el marco de la
Guía para la Expresión de la Incertidumbre de Medición
(GUM, JCGM~100:2008), considerando las fuentes de incertidumbre
relevantes para la medición de potencia activa.

\subsubsection{Fuentes de incertidumbre identificadas}

\begin{itemize}
    \item \textbf{Tipo~A — dispersión estadística de las mediciones:}
          incertidumbre debida a la variabilidad aleatoria del proceso
          de medición, cuantificada por la desviación estándar de la
          media $u_{\bar{P}}$ (ecuación~\ref{eq:incertidumbre_estadistica}).

    \item \textbf{Tipo~B — exactitud del sensor PZEM-004T:}
          la hoja de datos del módulo especifica una exactitud de
          $\pm$0{,}5\% en potencia activa.
          Asumiendo distribución rectangular, la incertidumbre
          estándar asociada es:
          \begin{equation}
              u_B = \frac{0{,}005 \times \bar{P}}{\sqrt{3}}
                  = \PLACEHOLDER{X{,}X}~\text{W}
              \label{eq:incertidumbre_B}
          \end{equation}

    \item \textbf{Tipo~B — variación de la tensión de red:}
          fluctuaciones de la tensión de alimentación durante el
          ensayo se propagan a la medición de potencia.
          Con una variación observada de $\Delta V = \PLACEHOLDER{X{,}X}$~V
          y el coeficiente de sensibilidad
          $c_V = \partial P / \partial V \approx I_{\text{rms}}$:
          \begin{equation}
              u_{V} = c_V \cdot \frac{\Delta V}{\sqrt{3}}
                    = \PLACEHOLDER{X{,}X}~\text{W}
              \label{eq:incertidumbre_V}
          \end{equation}

    \item \textbf{Tipo~B — resolución del sensor:}
          la resolución en potencia del PZEM-004T es 0{,}1~W.
          Asumiendo distribución rectangular:
          \begin{equation}
              u_{\text{res}} = \frac{0{,}1}{2\sqrt{3}}
                             \approx 0{,}029~\text{W}
              \label{eq:incertidumbre_res}
          \end{equation}
\end{itemize}

\subsubsection{Incertidumbre combinada y expandida}

Considerando las fuentes anteriores como independientes entre sí,
la incertidumbre estándar combinada es:

\begin{equation}
    u_c = \sqrt{u_{\bar{P}}^2 + u_B^2 + u_V^2 + u_{\text{res}}^2}
        = \PLACEHOLDER{X{,}X}~\text{W}
    \label{eq:incertidumbre_combinada}
\end{equation}

La incertidumbre expandida, con factor de cobertura $k=2$
(nivel de confianza $\approx$95\%, distribución normal), es:

\begin{equation}
    U = k \cdot u_c = 2 \times \PLACEHOLDER{X{,}X}
      = \PLACEHOLDER{X{,}X}~\text{W}
    \label{eq:incertidumbre_expandida}
\end{equation}

El resultado final de la medición de potencia se expresa como:

\begin{equation}
    P = (\PLACEHOLDER{XXXX{,}X} \pm \PLACEHOLDER{X{,}X})~\text{W}
        \quad (k=2,\; p \approx 95\%)
    \label{eq:resultado_final}
\end{equation}

\subsubsection{Relación entre incertidumbre y límite normativo}

La incertidumbre expandida del sistema ($U = \PLACEHOLDER{X{,}X}$~W)
representa el \PLACEHOLDER{X{,}X}\% del límite normativo
($\pm$10\% de $P_{\text{dec}} = \PLACEHOLDER{XXXX}$~W).
Este valor es \PLACEHOLDER{aceptable para los propósitos del
prototipo / recomendable reducir mediante calibración del sensor},
dado que el criterio general de metrología establece que la
incertidumbre del instrumento de medición debe ser al menos
cuatro veces menor que la tolerancia que se desea verificar.

% ------------------------------------------------------------
\subsection{Evaluación de la repetibilidad del sistema}
\label{subsec:repetibilidad}
% ------------------------------------------------------------

La repetibilidad del sistema se evalúa como la capacidad de
obtener resultados consistentes bajo las mismas condiciones
experimentales en mediciones sucesivas.

\begin{table}[H]
\centering
\caption{Indicadores de repetibilidad del ensayo de potencia}
\label{tab:repetibilidad}
\renewcommand{\arraystretch}{1.35}
\begin{tabular}{@{} l r l @{}}
\toprule
\textbf{Indicador} & \textbf{Valor} & \textbf{Interpretación} \\
\midrule
Media $\bar{P}$                 & \PLACEHOLDER{XXXX,X} W  & Valor central de la distribución \\
Desviación estándar $s_P$       & \PLACEHOLDER{X,X} W     & Dispersión de las mediciones \\
Coeficiente de variación CV     & \PLACEHOLDER{X,X}\%     & Repetibilidad relativa \\
Rango ($P_{\max} - P_{\min}$)   & \PLACEHOLDER{X,X} W     & Amplitud total de dispersión \\
Incertidumbre estándar $u_{\bar{P}}$ & \PLACEHOLDER{X,X} W & Incertidumbre tipo~A \\
\midrule
\textbf{Criterio de repetibilidad} & CV $<$ 2\%
    & \PLACEHOLDER{\textcolor{green!60!black}{\textbf{CUMPLE}} /
                   \textcolor{red}{\textbf{NO CUMPLE}}} \\
\bottomrule
\end{tabular}
\end{table}

% ------------------------------------------------------------
\subsection{Discusión de resultados}
\label{subsec:discusion}
% ------------------------------------------------------------

Los resultados obtenidos permiten extraer las siguientes
conclusiones sobre el desempeño del prototipo:

\paragraph{Exactitud en la medición de potencia.}
El error relativo promedio frente al instrumento de referencia
fue de \PLACEHOLDER{X{,}XX}\%, valor que
\PLACEHOLDER{se encuentra dentro / supera levemente}
la especificación del fabricante del módulo PZEM-004T (±0{,}5\%).
\PLACEHOLDER{Describir aquí la causa probable del error observado
si supera la especificación, p.ej.: caída de tensión en el cableado,
temperatura de operación del módulo, saturación del TC, etc.}

\paragraph{Repetibilidad.}
El coeficiente de variación obtenido de \PLACEHOLDER{X{,}X}\%
\PLACEHOLDER{confirma / no confirma} una repetibilidad adecuada
del sistema para el propósito del ensayo.
\PLACEHOLDER{Si CV > 2\%: identificar causa de la variabilidad
(fluctuaciones de red, termostato del DUT activo durante la medición,
interferencias en el bus UART, etc.) y proponer acción correctiva.}

\paragraph{Cumplimiento normativo.}
El DUT ensayado \PLACEHOLDER{aprobó / reprobó} el criterio de
consumo de la norma IEC~60335-1 (cláusula~10.2), con una desviación
de \PLACEHOLDER{X{,}X}\% respecto al valor declarado,
dentro del límite admisible de $\pm$10\%.
Este resultado \PLACEHOLDER{es / no es} consistente con lo esperado
para un electrodoméstico \PLACEHOLDER{en buenas condiciones / con
el defecto identificado: descripción}.

\paragraph{Limitaciones identificadas.}
Durante la campaña de ensayos se identificaron las siguientes
limitaciones del prototipo en su estado actual:

\begin{itemize}
    \item \PLACEHOLDER{Limitación 1: p.ej., el módulo PZEM-004T
          presenta lecturas inestables durante los primeros 30~s
          de conexión, por lo que se implementó un retardo de
          estabilización en el software.}
    \item \PLACEHOLDER{Limitación 2: p.ej., la precisión de la
          medición de factor de potencia se degrada para cargas
          muy resistivas ($\cos\varphi > 0{,}98$) debido a la
          resolución de 0{,}01 del módulo.}
    \item \PLACEHOLDER{Limitación 3: agregar según observaciones
          reales del ensayo.}
\end{itemize}

Estas limitaciones no comprometen la validez del ensayo de consumo
como función principal del prototipo, pero deben ser consideradas
al momento de ampliar el sistema hacia los ensayos de fuga y
rigidez dieléctrica contemplados en la arquitectura del equipo.

% ------------------------------------------------------------
\subsection{Comparación con el criterio de éxito del proyecto}
\label{subsec:criterio_exito}
% ------------------------------------------------------------

La Tabla~\ref{tab:criterios_exito} evalúa los resultados obtenidos
frente a los criterios de éxito definidos en la
Sección~\ref{sec:criterios_evaluacion} del presente informe.

\begin{table}[H]
\centering
\caption{Evaluación de criterios de éxito del prototipo}
\label{tab:criterios_exito}
\renewcommand{\arraystretch}{1.4}
\begin{tabular}{@{} p{5cm} p{3cm} p{3cm} p{1.8cm} @{}}
\toprule
\textbf{Criterio} & \textbf{Meta definida} & \textbf{Resultado obtenido} & \textbf{Estado} \\
\midrule
Reducción del tiempo de prueba vs.\ método manual
    & $\geq 30\%$
    & \PLACEHOLDER{XX\%}
    & \PLACEHOLDER{\textcolor{green!60!black}{\checkmark} /
                   \textcolor{red}{$\times$}} \\
\addlinespace
Reducción de intervención del operario
    & $\geq 80\%$
    & \PLACEHOLDER{XX\%}
    & \PLACEHOLDER{\textcolor{green!60!black}{\checkmark} /
                   \textcolor{red}{$\times$}} \\
\addlinespace
Error relativo en potencia frente a patrón
    & Según especif.\ PZEM ($\pm$0{,}5\%)
    & \PLACEHOLDER{$\pm$X{,}XX\%}
    & \PLACEHOLDER{\textcolor{green!60!black}{\checkmark} /
                   \textcolor{red}{$\times$}} \\
\addlinespace
Cumplimiento normativo (criterio IEC~10.2)
    & Detección correcta APROBADO / REPROBADO
    & \PLACEHOLDER{Correcto / Incorrecto}
    & \PLACEHOLDER{\textcolor{green!60!black}{\checkmark} /
                   \textcolor{red}{$\times$}} \\
\addlinespace
Repetibilidad (CV)
    & CV $< 2\%$
    & \PLACEHOLDER{X{,}X\%}
    & \PLACEHOLDER{\textcolor{green!60!black}{\checkmark} /
                   \textcolor{red}{$\times$}} \\
\addlinespace
Ajuste al presupuesto
    & $\leq$ USD~1030
    & USD~\PLACEHOLDER{XXX}
    & \PLACEHOLDER{\textcolor{green!60!black}{\checkmark} /
                   \textcolor{red}{$\times$}} \\
\bottomrule
\end{tabular}
\end{table}
